\appendix

\section{FAQs}
\begin{enumerate}
    \item \textbf{Dataset or Benchmark:} Is this a dataset or a benchmark? {\textbf{A benchmark}}
    
    \item \textbf{Benchmark:} For benchmarks, the supplementary materials must ensure that all results are easily reproducible (i.e., all necessary datasets, code, and evaluation procedures must be accessible and documented)
    % ~\\
    
    \textbf{Datasets:} \discoverybench{} is released at \url{https://github.com/allenai/discoverybench/tree/main/discoverybench} . 

    \textbf{Code (Baseline Models):} Code for Discovery Agents are provided in the repository, at: \url{https://github.com/allenai/discoverybench/tree/main/agents}. A CLI is available to run the discovery agents on the benchmark. 

    \textbf{Evaluation Procedures:} Please follow our main paper for the details of our evaluation process. The code to run eval on a single instance of our benchmark is provided at: \url{https://github.com/allenai/discoverybench/tree/main/eval}. A CLI and some example scripts have been provided as well. 

    \item{\textbf{Accessibility}: The following are accessibility items on the submission checklist:}

    \textbf{Links to access the benchmark:} The link to access the benchmark is provided in the main submission (\url{https://github.com/allenai/discoverybench/tree/main/discoverybench}). 

    \textbf{Any data should use open and widely used formats. Simulation environments should explain how they can be used:} Our data are stored in widely accessible standard formats (e.g., \textsc{JSON, CSV}), with the structure described in  \Cref{sec:datadetails}.  

    \textbf{Long-term preservation.} Code and data are provided on \textsc{Github}.  All aspects will be publicly available for a long term.

    \textbf{Explicit Licence:} Our benchmark is licensed using \textsc{ODC-BY} and the associated code is licensed with \textsc{Apache 2.0}, as included in the \textsc{Github} repository. 

    \textbf{Structured Metadata for a dataset:} Our dataset is also available as the HuggingFace dataset: \url{https://huggingface.co/datasets/allenai/discoverybench}. Structured Metadata will be available once we finalize our work after addressing the reviewers' comments, if any.
    % \bodhi{TBA} This is not applicable (\textsc{DiscoveryWorld} is a benchmark that takes place in a virtual environment, and there is no typical data release as such).  The associated data from our baseline runs is included as links to Google Drive folders on the repository (\url{https://github.com/allenai/discoveryworld/tree/main/data}).

    \textbf{A persistent dereferenceable identifier (e.g., a code repository such as GitHub):} The repository for our benchmark is: \url{https://github.com/allenai/discoverybench}.

\end{enumerate}

\section{Limitations}
\label{sec:limitations}
We currently filtered domains and tasks that required forecasting, simulation, or very specific modeling (species distribution, infection spread, astrophysics equations for exoplanets) in the benchmark as they were very time-consuming to replicate as well as discover hypotheses. As a result, we discarded more papers focused on natural and physical sciences compared to social sciences, which we plan to include in future benchmarks. 

We currently do not tackle the challenge of understanding and processing massive datasets, such as the 8.92 petabytes data from the Cancer Genome Atlas (\url{https://portal.gdc.cancer.gov}) or the extensive brain data from the Allen Institute (\url{https://alleninstitute.org/division/brain-science}). While the potential to discover new insights from such vast data volumes is significant, ensuring these findings are robust and not subject to $p$-hacking remains unaddressed by our current methods.

We currently do not handle multi-modal data and complex pipelines, such as those needed for analyzing satellite and other geospatial data relevant to climate science and astronomy data. This would involve multiple stages of data processing, the use of various tools, and managing workflow complexities, for example, analyzing thousands of species patterns combined with satellite data to study habitats. So we do not incorporate workflows like those of EarthRanger (\url{https://www.earthranger.com}).


\newpara{Ethical Considerations} There could be many potential societal consequences of systems tuned on our proposed benchmark since it involves using LLMs, such as policy misuse, legal ramifications, and false discovery. On the positive side, our proposed benchmark can advance the rate of discovery, leading to an improved standard of living and social well-being.


\section{Data collection for \real{}}
\label{sec:real}
For data-first approach, replication took 15 to 40 person-hours for each NLS-related paper and up to 90 person-hours
for the GBIF dataset, where specialized domain knowledge and tools led to higher complexity. All papers replicated in the NLS dataset were included, while less than half of the papers in specialized datasets like GBIF and WBOD were added to \discoverybench{}. 

\textbf{Citation/Repositories for \real{}:}
List of scientific works from where we have replicated our gold workflows and hypotheses:
\begin{enumerate}
    \item Sociology: \cite{zaw2016race,apel2010impact,alexander1982social, smith2005time,dougherty2003numeracy}
    \item Biology: \cite{cerezer2023accelerated, riera2024effect}
    \item Economics: \cite{pal2023impact,appiah2017effect,weatherly2022impact,rambeli2021dynamic,ottaviano2013immigration}
    \item Engineering: \cite{alves2023status}
    \item Meta-science: \cite{heyard2024meta}
    \item Humanities: \cite{brozio2024patterns,brozio2019monuments,lorenz2018kommunikationsstrukturen,maida2023unter,sommerfeld2013gerategeld,feeser2019human,parkinson2021radiocarbon,kneisel2021chronology,kaiser2021time,brinkmann2019copper,dal2019s,palmisano2021long,parkinson2021radiocarbon}
\end{enumerate}
All assets come under CC license or open licenses. 


% \begin{itemize}
%     \item Race, Wealth and Incarceration: Results from the National Longitudinal Survey of Youth
%     \item The impact of incarceration on employment during the transition to adulthood
%     \item Social Background, Academic Resources, and College Graduation: Recent Evidence from the National Longitudinal Survey
%     \item Are time preference and body mass index associated?: Evidence from the National Longitudinal Survey of Youth
%     \item Numeracy, literacy and earnings: evidence from the National Longitudinal Survey of Youth
%     \item Accelerated body size evolution in upland environments is correlated with recent speciation in South American freshwater fishes
%     \item Effect of introduction pathways on the invasion success of non-native plants along environmental gradients
%     \item Pal, L.C. (2023). Impact of Education on Economic Development. Khazanah Pendidikan Islam.
%     \item The Effect of Education Expenditure on Per Capita GDP in Developing Countries
%     \item Weatherly, H., Lopez, K., \& Tierra, C.J. (2022). The Impact of Education on GDP per Capita.
%     \item Norimah Rambeli, Dayang Affizzah Awang Marikan, Jan M. Podivinsky, Rosilawati Amiruddin, & Ismadi Ismail. (2021). The Dynamic Impact of Government Expenditure in Education on Economic Growth. International Journal of Business and Society, 22(3), 1487–1507. https://doi.org/10.33736/ijbs.4318.2021
%     \item Industrial Practices of Requirements Engineering for ML-Enabled Systems in Brazil
%     \item Heyard, R. and Held, L. (2024). Meta-regression to explain shrinkage and heterogeneity in large-scale replication projects
%     \item Patterns of Socio-economic Cultural Transformations in Neolithic and Bronze Age Societies in the Central Northern European Plain
%     \item Brozio, J. P., Müller, J., Furholt, M., Kirleis, W., Dreibrodt, S., Feeser, I., Dörfler, W., Weinelt, M., Raese, H., \& Bock, A. (2019). Monuments and economies: What drove their variability in the middle-Holocene Neolithic? The Holocene [Special Issue: Scales of Transformation – Human-Environmental Interaction in Prehistoric and Archaic Societies], 29(10), 1558–1571. \url{https://doi.org/10.1177/0959683619857227}
%     \item Lorenz, L. (2018). Kommunikationsstrukturen mittelneolithischer Gesellschaften im nordmitteleuropäischen Tiefland. Dr. Rudolf Habelt.
%     \item Ahrens, C. (1966). Vorgeschichte des Kreises Pinneberg und der Insel Helgoland. Wachholtz.
%     \item Schaefer-Di Maida, S. (2023). Unter Hügeln. Bronzezeitliche Transformationsprozesse in Schleswig-Holstein am Beispiel des Fundplatzes von Mang de Bargen (Bornhöved, Kr. Segeberg) [Doctoral dissertation, Christian-Albrechts-Universität zu Kiel]. Sidestone Press. \url{https://doi.org/10.59641/q9013tc}
%     \item Sommerfeld, C. (1994). Gerätegeld Sichel. Studien zur monetären Struktur bronzezeitlicher Horte im nördlichen Mitteleuropa. De Gruyter.
%     \item Feeser, I., Dörfler, W., Kneisel, J., Hinz, M., \& Dreibrodt, S. (2019). Human impact and population dynamics in the Neolithic and bronze age: Multi-proxy evidence from north-western Central Europe. The Holocene [Special Issue: Scales of Transformation – Human-Environmental Interaction in Prehistoric and Archaic Societies], 29(10), 1596–1606. \url{https://doi.org/10.1177/0959683619857223}
%     \item Parkinson, E. W., McLaughlin, T. R., Esposito, C., \& Stoddart, \& Malone, C. (2021). Radiocarbon dated trends and Central Mediterranean prehistory. Journal of World Prehistory, 34, 317–379. \url{https://doi.org/10.1007/s10963-021-09158-4}
%     \item Kneisel, J. (2021). Chronology and transformation. The transition from Bronze to Iron Age in Northern Europe. In E. Kaiser \& W. Schier (Eds.), Time and materiality: Periodization and regional chronologies at the transition from Bronze to Iron Age in Eurasia (1200–600 BCE) (pp. 237–263). VML Verlag Marie Leidorf.
%     \item Brinkmann, J. (2019). Copper output, demand for wood and energy expenditure – Evaluating economic aspects of Bronze Age metallurgy. In M. Dal Corso, W. Kirleis, J. Kneisel, N. Taylor, M. Wieckowska-Lüth, \& M. Zanon (Eds.), How’s life? Living conditions in the 2nd and 1st Millennia BCE (pp. 11–34). Sidestone Press. \url{https://doi.org/10.59641/e7245hq}
%     \item Palmisano et al. (2021) “Long-Term Demographic Trends in Prehistoric Italy: Climate Impacts and Regionalised Socio-Ecological Trajectories” (DOI: \url{https://doi.org/10.1007/s10963-021-09159-3})
%     \item Parkinson et al. (2021) “Radiocarbon Dated Trends and Central Mediterranean Prehistory” (DOI: \url{https://doi.org/10.1007/s10964-021-09158-4})
%     \item Gianmarco I. P. Ottaviano \& Giovanni Peri \& Greg C. Wright, 2013. "Immigration, Offshoring, and American Jobs," American Economic Review, American Economic Association, vol. 103(5), pages 1925-59, August.
% \end{itemize}




\section{Data Generation for \synth{}}
\label{sec:synth}
For leaves, we use different sampling strategies based on the data type. 
Specifically, for categorical nodes, we sample instances with replacement from the range of allowed values, whereas for numeric, we first select a distribution (e.g., normal) and its parameters based on the specified range and then perform sampling. 
For each subsequent level in $\mathcal{F}$, we create new columns for nodes by simply executing their pandas expressions\footnote{The expression is guaranteed to only have variables already generated due to the bottom-up construction.}. 
To recover from any execution errors, we additionally use a self-refine \citep{madaan2023selfrefine} approach to generate new pandas expressions guided by the execution error logs. 
Finally, to mimic real-world challenges in data collection, we probabilistically perturb each instance $x \in \mathbf{x}_i$ 
% with $0.3$ 
% \tushar{seems high. Can it be grounded with any of the real benchmarks?} \dhruv{This would be nice but we don't have support for this in the real-data since the two sets were developed in parallel.} 
% probability 
by adding noise or dropping values 
% \tushar{would you lose 15\% of the data?} \dhruv{The implementation is that every \textit{cell} in the dataset has a 30\% chance of being perturbed, so we do lose/noise 30\% of cells on average in the full dataset, but not entire rows.} 
to create missing data\footnote{Each value is noised independently; therefore, each row has sufficient true data useful for discovery.}. 
After generation, $D$ contains a column for each node in $\mathcal{F}$.


\section{Datasheets}
\label{sec:datasheets}

\subsection{Motivation}

\begin{itemize}
    \item \textbf{For what purpose was the dataset created?} \discoverybench{} is created to help assess large language models' (LLMs) ability to automate the search and verification of hypotheses purely from a set of provided datasets.
    \item \textbf{Who created the dataset (e.g., which team, research group) and on behalf of which entity (e.g., company, institution, organization)?} Authors belong to the Allen Institute for AI, OpenLocus, and the University of Massachusetts Amherst. The data collection is part of research efforts conducted by the Allen Institute for AI.
    \item \textbf{Who funded the creation of the dataset?} Allen Institute for AI.
\end{itemize}

% \subsection{Composition}

% \begin{itemize}
%     \item \textbf{What do the instances that comprise the dataset represent?}
%     \item \textbf{How many instances are there in total (of each type, if appropriate)?}
%     \item \textbf{Is there a label or target associated with each instance?}
%     \item \textbf{Are there recommended data splits (e.g., training, development/validation, testing)?}
% \end{itemize}

\subsection{Collection Process}

\begin{itemize}
    \item \textbf{How was the data associated with each instance acquired?}
    Our goal is to replicate the scientific process undertaken by researchers to search for and validate a hypothesis from one or more datasets. We focus on six scientific domains where data-driven research is the cornerstone of scientific progress: sociology, biology, humanities, economics, engineering, and meta-science. Our data collection follows either a \textbf{data-first} or \textbf{code-first} approach. Each instance has been manually implemented and verified by the authors for solvability.
\end{itemize}

\subsection{Uses}
\begin{itemize}
    \item \textbf{Has the dataset been used for any tasks already?}
    We use this benchmark to evaluate LLM's ability to search and verify hypotheses purely from a set of datasets.
    \item \textbf{Are there tasks for which the dataset should not be used?} We do not expect the community members to use this data to train models that can aggravate $p$-hacking.
\end{itemize}

\subsection{Distribution and Maintainance}
\begin{itemize}
    \item \textbf{How will the dataset will be distributed?}
    We distribute this benchmark via our \textsc{Github} repository: \url{https://github.com/allenai/discoverybench} and \url{https://huggingface.co/datasets/allenai/discoverybench}.
    \item \textbf{How can the owner/curator/manager of the dataset be contacted?} For any benchmark-related queries, please contact: \texttt{bodhisattwam@allenai.org}. For any code-related discussions, please raise an issue in \textsc{Github}: \url{https://github.com/allenai/discoverybench}.
\end{itemize}

\section{Composition of \discoverybench{}}
\label{sec:datadetails}

% Cover things like:
% 1) folder structure
% 2) data file structure - explain how 1 file have more than 1 queries

% \subsection*{Real Dataset Structure}
% TODO: BUG IN ENV THAT NEEDS FIXING
% \begin{forest}
%   for tree={
%     font=\ttfamily,
%     grow'=0,
%     child anchor=west,
%     parent anchor=south,
%     anchor=west,
%     calign=first,
%     edge path={
%       \noexpand\path [draw, \forestoption{edge}]
%       (!u.south west) + (0.2cm,0) |- (.child anchor)\forestoption{edge label};
%     },
%     before typesetting nodes={
%       if n=1
%         {insert before={[,phantom]}}
%         {}
%     },
%     fit=band,
%     before computing xy={l=15pt},
%   }
% [/
%   [test
%     [archaeology]
%     [introduction\_pathways\_non-native\_plants]
%     [meta\_regression]
%     [meta\_regression\_raw]
%     [nls\_incarceration]
%     [nls\_raw]
%     [nls\_ses]
%     [requirements\_engineering\_for\_ML\_enabled\_systems]
%     [worldbank\_education\_gdp
%       [worldbank\_education\_gdp\_indicators]
%     ]
%   ]
%   [train
%     [evolution\_freshwater\_fish]
%     [immigration\_offshoring\_effect\_on\_employment]
%     [nls\_bmi]
%     [nls\_bmi\_raw]
%   ]
% ]
% \end{forest}

\subsection{Metadata structure}
\begin{itemize}
    \item \textbf{id}: An identifier for the metadata.
    
    \item \textbf{domain}: The broad field of study or area of research.
    
    \item \textbf{workflow\_tags}: A set of keywords summarizing the main processes or techniques used in the replication implementation. They provide an overview of the methodological approach and facilitating the identification of relevant analytical techniques.
    
    \item \textbf{domain\_knowledge}: 
    \begin{itemize}
        \item Contextual information or insights related to the dataset, explaining how certain behaviors or variables can be interpreted within the field of study.
        \item It helps open avenues to think in directions that LLM might not have considered otherwise, broadening the understanding of the field.
    \end{itemize}
    
    \item \textbf{datasets}: Contains detailed information about the datasets used, including:
    \begin{itemize}
        \item \textbf{name}: The name or filename of the dataset.
        
        \item \textbf{description}: A summary of the dataset's contents and the type of data it includes.
        
        \item \textbf{max\_depth}: The maximum hierarchical level of nested data structures within the dataset, indicating the complexity of the data.
        
        \item \textbf{columns}: Detailed descriptions of each column in the dataset, including:
        \begin{itemize}
            \item \textbf{name}: The column's name or header.
            \item \textbf{description}: Explanation of the data contained in the column and its significance.
            \item \textbf{depth}: The hierarchical level of the column within the dataset, indicating its structural position.
        \end{itemize}
    \end{itemize}
    
%    \item \textbf{intermediate}: Intermediate results or processed data that are part of the analysis but not the final output, serving as steps or components in the analytical workflow -TBA
    
    \item \textbf{hypotheses}: Statements or predictions being tested, divided into:
    \begin{itemize}
        \item \textbf{main}: Primary hypotheses that are central to the discovery task.
%        \item \textbf{intermediate}: Secondary or supporting hypotheses that contribute additional context or detail to the main hypotheses, helping to refine or elaborate on the primary findings.
    \end{itemize}
    
    \item \textbf{workflow}: A step-by-step description of the replication process followed to validate the hypotheses, outlining the methods and procedures used from data preparation to final analysis. Some of the workflows and sub-workflows are high-level and thus the same for different queries as they follow the same implementation leading to a range of hypotheses. 
    
    \item \textbf{queries}: Goals related to each hypothesis, each including:
    \begin{itemize}
        \item \textbf{qid}: A unique identifier for the goal for a given true/gold hypothesis.
        \item \textbf{difficulty}: Categorization of the difficulty. Structurally defined for \synth{} using the semantic tree definition.
        \item \textbf{true\_hypothesis}: The hypothesis being tested through the goal. This defines the primary statement or prediction under investigation.
%        \item \textbf{true\_hypothesis\_expr}:
            % The mathematical or logical expression of the true hypothesis. Only appears for \synth{}.
        \item \textbf{relevant\_cols}: Columns from the dataset that are relevant to answering the query, indicating the specific data points that can be used in the analysis. Only appears for \synth{}.
        \item \textbf{target\_col}: The column being predicted or the dependent variable in the analysis. Only appears for \synth{}.
%        \item \textbf{prior\_findings}: Any previous research findings related to the hypothesis. Only appears for \synth{}.
        \item \textbf{question\_type}: The type of question being asked categorizing the nature of the inquiry.
%        \item \textbf{question\_specificity}: The level of detail or specificity of the question.
        \item \textbf{question}: The discovery goal.
    \end{itemize}
\end{itemize}



\subsection{Directory structure for \real{}}

% \paragraph{The real dataset have all the aforementioned keys except for difficulty and expression related ones used for evaluating synth.}

There may be more than one query per metadata. The train split contains 14 metadata files and 25 queries. The test split contains 144 metadata files and 239 queries. Metadata folders with the same prefixes use the same underlying dataset with either a subset or a preprocessed version. When dealing with a full dataset (i.e., nls\_raw), the task becomes substantially harder due to the data preparation required.

\begin{verbatim}
    |-test
    |---archaeology
    |---introduction_pathways_non-native_plants
    |---meta_regression
    |---meta_regression_raw
    |---nls_incarceration
    |---nls_raw
    |---nls_ses
    |---requirements_engineering_for_ML_enabled_systems
    |---worldbank_education_gdp
    |---worldbank_education_gdp_indicators
    |-train
    |---evolution_freshwater_fish
    |---immigration_offshoring_effect_on_employment
    |---nls_bmi
    |---nls_bmi_raw
\end{verbatim}

\subsection{Directory structure for \synth{}}

% \paragraph{The real dataset have all the aforementioned keys except for workflow, workflow\_tags, domain\_knowledge.}

% \paragraph{
There is one query per metadata. The train split contains 551 metadata files (queries), the dev split contains 153 metadata files (queries), and the test split contains 200 metadata files (queries).

\begin{verbatim}
   |-test
   |---ancient-languages_*_*
   |---artificial-ecosystems_*_*
   |---astronomy_*_*
   |---board-games_*_*
   |---coding-competitions_*_*
   |---digital-artistry_*_*
   |---futuristic-technology_*_*
   |---impressionist-art_*_*
   |---machine-learning_*_*
   |---molecular-gastronomy_*_*
   |---neuroscience_*_*
   |---philosophical-debates_*_*
   |---robotics_*_*
   |-train
   |---adventure-travel_*_*
   |---ancient-architecture_*_*
   |---ancient-astronomy_*_*
   |---aviation_*_*
   |---biodiversity-conservation_*_*
   |---cryptic-puzzles_*_*
   |---cryptocurrency_*_*
   |---culinary-arts_*_*
   |---cybersecurity_*_*
   |---environmental-activism_*_*
   |---fashion-design_*_*
   |---fine-arts_*_*
   |---literary-classics_*_*
   |---marine-biology_*_*
   |---marine-conservation_*_*
   |---medieval-literature_*_*
   |---musical-therapy_*_*
   |---photography_*_*
   |---robotic-explorers_*_*
   |---solar-power_*_*
   |---space-tourism_*_*
   |---steampunk-culture_*_*
   |---theater-productions_*_*
   |---underwater-archaeology_*_*
   |---urban-gardening_*_*
   |---vintage-automobiles_*_*
   |---virtual-reality_*_*

\end{verbatim}

\section{Discovery Agent}

% Cover how discovery agents could be run:

The command \texttt{discovery\_agent.py} is used with various options to customize its behavior for discovery tasks. Below are the options explained:

\begin{itemize}[leftmargin=*]
    \item \textbf{Usage:} \texttt{discovery\_agent.py [OPTIONS] QUERY} -- Executes the discovery agent with specified options.

    \item \textbf{Options:}
    \begin{itemize}
        \item \texttt{--agent\_type [coder|react]}: Specifies the type of agent to use for discovery. The default type is \texttt{coder}. Options include \texttt{coder} for code-related tasks and \texttt{react} for reactive tasks.
        
        \item \texttt{--model\_name TEXT}: Sets the model to be used. The default is \texttt{gpt-4o}. Available models include \texttt{gpt-4-turbo}, \texttt{llama-3-70b-chat}, \texttt{claude-3-opus}, and \texttt{gemini-pro}. An exhaustive list is available in \texttt{config/model\_config.json}.
        
        \item \texttt{--api\_config TEXT}: Path to the API configuration file. The default path is \texttt{config/api\_config.json}.
        
        \item \texttt{--log\_file TEXT}: Specifies the path to the log file where operations details are stored.
        
        \item \texttt{--metadata\_path TEXT}: Path to the metadata file. This option is required.
        
        \item \texttt{--metadata\_type [real|synth]}: Specifies the type of metadata, where \texttt{real} stands for actual metadata and \texttt{synth} for synthetic. This option is required.
        
        \item \texttt{--add\_domain\_knowledge}: Includes domain-specific knowledge in the query processing.
        
        \item \texttt{--add\_workflow\_tags}: Includes workflow tags in the query to enhance context.
        
        \item \texttt{--help}: Displays the help message and exits, showing all available command options.
    \end{itemize}
\end{itemize}

\section{Evaluation}

Explain about evaluation in a line and then explain the CLI usage here. 

The command \texttt{discovery\_eval.py} is used to evaluate the outputs generated by the discovery agent. Below are the detailed descriptions of the command options:

\begin{itemize}[leftmargin=*]
    \item \textbf{Usage:} \texttt{discovery\_eval.py [OPTIONS] QUERY} -- Executes the evaluation agent with specified options and a query.

    \item \textbf{Options:}
    \begin{itemize}
        \item \texttt{--gold\_hypo TEXT}: Specifies the gold standard hypothesis for comparison. This field is required.
        
        \item \texttt{--gold\_workflow TEXT}: Specifies the gold standard workflow to be used as a reference during evaluation.
        
        \item \texttt{--pred\_hypo TEXT}: Specifies the predicted hypothesis generated by the discovery agent. This field is required.
        
        \item \texttt{--pred\_workflow TEXT}: Specifies the predicted workflow generated by the discovery agent.
        
        \item \texttt{--metadata\_path TEXT}: Specifies the path to the metadata file that is utilized during evaluation. This field is required.
        
        \item \texttt{--metadata\_type [real|synth]}: Determines the type of metadata used in the evaluation, where \texttt{real} indicates actual metadata and \texttt{synth} indicates synthetic metadata. This field is required.
        
        \item \texttt{--eval\_output\_path TEXT}: Specifies where the evaluation results should be saved.
        
        \item \texttt{--help}: Displays the help message and exits, detailing all available command options.
    \end{itemize}
\end{itemize}


\section{Experiments}
\label{app:exp}
For GPT-based models, we use OpenAI API (\url{https://platform.openai.com/docs/models}), and for Llama3, we used Together API (\url{https://docs.together.ai/docs/inference-models})

\section{Evaluator Prompts}
\label{app:evaluation-prompts}

We provide below the exact prompts used for our GPT-4 based evaluation of the generated hypothesis against the gold hypothesis.

\begin{listing}[!ht]
\begin{minted}[frame=lines,
baselinestretch=1,
bgcolor=Box3Color,
fontsize=\footnotesize, breaklines, breaksymbolleft={}, breaksymbolright={}]{python}
decomposition_prompt = f"""\
        Given a set of dataset columns, a ground-truth hypothesis, and the analysis workflow used, your task is to extract the set of sub-hypotheses that are present in the hypothesis such that each sub-hypothesis covers a separate context, is self-sufficient, and operates on a coherent set of 3 dimensions: Context, Variables, and Relations.
        
        Here are the definitions for these dimensions:
        
        - Contexts: Boundary conditions that limit the scope of a sub-hypothesis. E.g., “for men over the age of 30”, “in Asia and Europe”, or "None" if there is no boundary condition specified.
        
        - Variables: Known concepts that interact in a meaningful way under a given context to produce the sub-hypothesis. E.g., gender, age, income, or "None" if there is no interacting variable.
        
        - Relations: Interactions between a given set of variables under a given context to produce the sub-hypothesis. E.g., “quadratic relationship”, “inversely proportional”, piecewise conditionals, or "None" if there is no interacting relationship. 
        
        Make sure to only use the information present in the hypothesis and the workflow. Do not add any new information.
        If no sub-hypotheses can be extracted, return an empty list.

        Here is the metadata for the task:
        ```json
        {{
            "datasets": {dataset_metadata},
            "hypothesis": "{hypothesis}",
            "workflow": "{workflow}"
        }}
        ```

        Return your answer as a JSON object in the following format:
        ```json
        {{
        "sub_hypo": [
            {{
                "text": the sub-hypothesis in natural language,
                "context": a short text description of the context of the sub-hypothesis,
                "variables": a list of columns involved in the sub-hypothesis,
                "relations": a short text description of the relationship between the variables of the sub-hypothesis,
                "explanation": a short text explanation for the breakdown of the sub-hypothesis
            }},
            ...
        ]
        }}```
        """
\end{minted}
\vspace{-1em}
\caption{Decomposition Prompt to obtain sub-hypotheses from a hypothesis.}
\label{list:decompose}
\end{listing}

\begin{listing}[!ht]
\begin{minted}[frame=lines,
baselinestretch=1,
bgcolor=Box3Color,
fontsize=\footnotesize, breaklines, breaksymbolleft={}, breaksymbolright={}]{python}
matching_prompt = f"""
        Given a gold hypothesis, a gold context, a predicted hypothesis, and a predicted context, your task is
        to determine if the predicted context semantically matches the ground-truth context.
        
        Here is the definition for Context: Boundary conditions that limit the scope of a sub-hypothesis. E.g., “for men over the age of 30”, “in Asia and Europe”, or "None" if there is no boundary condition specified.
        
        If the predicted context matches the gold context, return true, otherwise return false.

        Here is the metadata for the task:
        ```json
        {{
            "gold_hypothesis": "{gold_hypotheis}",
            "gold_context": "{gold_context}",
            "predicted_hypothesis": "{pred_hypothesis}",
            "predicted_context": "{pred_context}"
        }}
        ```

        Return your answer as a JSON object in the following format:
        ```json
        {{
            "match": true or false
        }}
        ```"""
\end{minted}
\vspace{-1em}
\caption{Matching prompt to match contexts of two sub-hypotheses.}
\label{list:match}
\end{listing}

\begin{listing}[!ht]
\begin{minted}[frame=lines,
baselinestretch=1,
bgcolor=Box3Color,
fontsize=\footnotesize, breaklines, breaksymbolleft={}, breaksymbolright={}]{python}
main_context = f"""
        You are going to compare two natural-language hypotheses HypoA and HypoB accompanied with optional workflows: WorkflowA for HypoA and WorkflowB for HypoB.
        Both the hypotheses answer the natural language query "QUERY" over the dataset(s) described by dataset description(s) and column description(s) below.
        Compare HypoA and HypoB in terms of three aspects: Contexts, Variables, and Relations.
        E.g., for the hypothesis "From 1995 to 2009, the number of sandhill cranes around the tundra (Indigilka River) surged by an astounding ~10X":
        * Contexts refer to the stratification of the data under which the given hypothesis is True. E.g., "For all women", "From 1995 to 2009". 
        * Variables refer to the set of variables (either dependent or independent) that are mentioned in the hypothesis. E.g., number of sandhill cranes, location.
        * Relations refer to the form of relation between the variables. E.g., "surged by ~10x".

        Answer the following questions for a given pair of hypotheses, HypoA and HypoB, along with an explanation grounded on the QUERY and the DATASET(S).

        Here is the metadata for the task:
        ```json
        {{
        "datasets": {datasets_json},
        "query": {query},
        "HypoA": {gold_hypo},
        "WorkflowA": {gold_workflow},
        "HypoB": {gen_hypo},
        "WorkflowB": {gen_workflow}
        }}
        ```

        {variable_question}"""
variable_question = """\
        Question: For both HypoA and HypoB, what are the different variables found in the hypotheses? \
        Return your answer as a JSON object in the following format:
        ```json
        {{
        "sizeA": num of variables used in HypoA
        "sizeB": num of variables used in HypoB
        "intersection": num of variables common in HypoA and HypoB. Use *fuzzy matching* to determine intersection, accounting for paraphrases or slightly different surface forms
        "explanation": a short text explanation about the variables
        }}```
        Answer:"""
\end{minted}
\vspace{-1em}
\caption{Prompt for variable alignment between two sub-hypotheses.}
\label{list:variable}
\end{listing}


\begin{listing}[!ht]
\begin{minted}[frame=lines,
baselinestretch=1,
bgcolor=Box3Color,
fontsize=\footnotesize, breaklines, breaksymbolleft={}, breaksymbolright={}]{python}
main_context = f"""
        You are going to compare two natural-language hypotheses HypoA and HypoB accompanied with optional workflows: WorkflowA for HypoA and WorkflowB for HypoB.
        Both the hypotheses answer the natural language query "QUERY" over the dataset(s) described by dataset description(s) and column description(s) below.
        Compare HypoA and HypoB in terms of three aspects: Contexts, Variables, and Relations.
        E.g., for the hypothesis "From 1995 to 2009, the number of sandhill cranes around the tundra (Indigilka River) surged by an astounding ~10X":
        * Contexts refer to the stratification of the data under which the given hypothesis is True. E.g., "For all women", "From 1995 to 2009". 
        * Variables refer to the set of variables (either dependent or independent) that are mentioned in the hypothesis. E.g., number of sandhill cranes, location.
        * Relations refer to the form of relation between the variables. E.g., "surged by ~10x".

        Answer the following questions for a given pair of hypotheses, HypoA and HypoB, along with an explanation grounded on the QUERY and the DATASET(S).

        Here is the metadata for the task:
        ```json
        {{
        "datasets": {datasets_json},
        "query": {query},
        "HypoA": {gold_hypo},
        "WorkflowA": {gold_workflow},
        "HypoB": {gen_hypo},
        "WorkflowB": {gen_workflow}
        }}
        ```

        {variable_question}"""
dimension_question = """
        Question: Does HypoB exhibit the same relation as HypoA?
        Compare using the following example hierarchy of relationships (based on specificity): \
        "there exists a relationship" > "positive relationship" > "positive AND (linear OR quadratic)" > "positive AND linear."
        Options: A) very similar B) similar but general than HypoA C) different
        Return your answer as a JSON object in the following format:
        ```json
        {{
        "answer": one of the options from A) very similar B) similar but general than HypoA C) different
        "explanation": a short text explanation about the relationship comparison
        }}```
        Answer:"""
\end{minted}
\vspace{-1em}
\caption{Prompt for relationship alignment between two sub-hypotheses.}
\label{list:rel}
\end{listing}


% \section{\real{} Task Examples}
% \label{app:real-examples}

% \begin{minted}[frame=lines,
% baselinestretch=1,
% bgcolor=Box3Color,
% fontsize=\footnotesize, breaklines, breaksymbolleft={}, breaksymbolright={}]{python}
% task = 
% f"""\
% Load all datasets using python using provided paths.
% Paths:
% ['../DiscoveryBench/archaeology/pollen_openness_score.csv', '../DiscoveryBench/archaeology/time_series_data.csv', '../DiscoveryBench/archaeology/capital.csv'].

% Dataset name: pollen_openness_score.csv
% Dataset description: Records of pollen data's PCA & interpolations across sites.
% Brief description of columns:
% Unnamed: Index or a time marker in years counting backward,
% calBP: Calibrated years Before Present (1950 AD),
% CE: Common Era,
% Belau_PC1: PC1 of principal components for pollen in Belau,
% Woserin_PC1: PC1 of principal components for pollen in Woserin,
% Belau_PC1_inter: Interpolated PC1 for the Belau site,
% Woserin_PC1_inter: Interpolated PC1 for the Woserin site,
% MEAN: The average of the interpolated PC1 for the Belau and Woserin sites,
% SMOOTH_MEAN_50y: Smoothed averages of the PC1 over 50 years,
% SMOOTH_MEAN_100y: Smoothed averages of the PC1 over 100 years, 
% SMOOTH_MEAN_150y: Smoothed averages of the PC1 over 150 years, 
% SMOOTH_MEAN_200y: Smoothed averages of the PC1 over 200 years, 
% SMOOTH_MEAN_250y: Smoothed averages of the PC1 over 250 years, 

% Dataset name: time_series_data.csv
% Dataset description: This dataset provides detailed quantification of archaeological findings over various time periods, measured in Z values for different cultural and economic indicators such as tools, house sizes, materials, and monument data.
% Brief description of columns:
% CE: Common Era (BCE x (-1)),
% calBP: Calibrated years before the present,
% kde_all_mean: Mean of kernel density estimation of all data points, 
% kde_all_std: Standard deviation of kernel density estimation of all data points,
% kde_all_detrend: KDE of data points after detrending,
% g_all_mean: Mean of KDE growth rates,
% g_all_std: Standard deviation of KDE growth rates,
% pollen: Pollen data of Belau Lake,
% pollen_inter: Interpolated and forward filled missing pollen values, pollen_detrend: Detrended pollen values from interpolated pollen values,
% pollen_inter_100: Rolling mean of the interpolated pollen data with a window size of 100,
% pollen_grate_100: Percentage change of interpolated pollen data, 
% HatchetSword: Z values for Hatchets and Swords,
% HatchetSword_inter: Interpolated z values for Hatchets and Swords, 
% Dagger: Z values for Daggers,
% Dagger_inter: Interpolated z value for Daggers,
% HouseSize: Z values for House Size in meter squared, 
% HouseSize_inter: Interpolated z values for House Sizes in meter squared,
% CopperGold: Z values for Copper and Gold,
% CopperGold_inter: Interpolated z values for Copper and Gold,
% Amber: Z values for Amber, Amber_inter: Interpolated z values for Amber,
% MonumentCount: Z values for Monument Count,
% MonumentCount_inter: Interpolated z values for Monument Count, 
% Depot: Z values for Depot,
% Depot_inter: Interpolated z values for Depot,
% Sickle: Z values for Sickle,
% Sickle_inter: Interpolated z values for Sickle,
% AxesCelts: Z values for Axes and Celts,
% AxesCelts_inter: Interpolated z values for Axes and Celts, 
% MonumentSize: Z values for Monument Size,
% MonumentSize_inter: Interpolated z values for Monument Size, 
% PotteryForm: Z values for Pottery Form,
% PotteryForm_inter: Interpolated z values for Pottery Form, 
% PotteryDecoration: Z values for Pottery Decoration, PotteryDecoration_inter: Interpolated z values for Pottery Decoration

% Dataset name: capital.csv
% Dataset description: This dataset contains archaeological data on various forms of capital across different prehistoric periods.
% Brief description of columns:
% BCE: Before Common Era,
% ZAxtSchwert: Z values for Hatchets and Swords,
% ZDolch: Z values for Daggers,
% Zhausgr: Z values for House Size,
% ZCU_AU: Z values for Copper and Gold,
% Zamber: Z values for Amber,
% ZMonument: Z values for Monument Count,
% ZHort: Z values for Depot,
% ZSichel: Z values for Sickle,
% ZBeil: Z values for Axes and Celts,
% ZMW: Z values for Monument Size,
% ZKeform: Z values for Pottery Form,
% Zkeverz: Z values for Pottery Decoration

% Goal: In which century did the axes become quantitatively most frequent?

% Target hypothesis: At the end of the 4th millennium BCE, axes become quantitatively most frequent.
% """
% \end{minted}

% \begin{minted}[frame=lines,
% baselinestretch=1,
% bgcolor=Box3Color,
% fontsize=\footnotesize, breaklines, breaksymbolleft={}, breaksymbolright={}]{python}
% task =
% f"""\
% Load all datasets using python using provided paths.
% Paths: ['../DiscoveryBench/nls_ses_processed/nls_ses_processed.csv'].

% Dataset name: nls_ses_processed.csv
% Dataset description: This dataset contains social background factors (race, gender, and socioeconomic status) and academic resources at the time of secondary school graduation (standardized test scores, class rank, and curriculum) for the participants of the NLS.
% Brief description of columns:
% CASE ID: Unique ID of each respondent,
% SAMPLE_RACE: Race of the respondent (Hispanic, Black or White), 
% SAMPLE_SEX: Sex of the respondent (Male or Female),
% FAMILY SIZE OF SAMPLE: Family size of the respondent,
% ABILITY: COMPOSITE OF ASVAB SCORE: Composite variable created by summing following variables: ASVAB - Arithmetic Reasoning Z Score (rounded), 1981, ASVAB - Word Knowledge Z Score (rounded), 1981, ASVAB - Paragraph Comprehension Z Score (rounded), 1981, ASVAB - Mathematics Knowledge Z Score (rounded), 1981,
% BA DEGREE COMPLETED: Boolean variable that equals 1 if the BA Degree was completed by the respondent else 0,
% PERCENTILE IN CLASS: Respondent's percentile in the class that he attended in school last year, SES: Socioeconomic Status of the respondent

% Goal: In which racial group was the strongest effect of Socioeconomic Status observed?

% Target hypothesis: The strongest effect of Socioeconomic Status was observed in the Black racial group.
% """
% \end{minted}


% \begin{minted}[frame=lines,
% baselinestretch=1,
% bgcolor=Box3Color,
% fontsize=\footnotesize, breaklines, breaksymbolleft={}, breaksymbolright={}]{python}
% task =
% f"""\
% Load all datasets using python using provided paths.
% Paths: ['../DiscoveryBench/nls_incarceration_processed/nls_incarceration_processed.csv'].

% Dataset name: nls_incarceration_processed.csv
% Dataset description: This dataset was created from the National Longitudinal Study of Youth's 1979 cohort to about the race, wealth, and incarceration status of the participants.
% Brief description of columns: 
% race: Race of the respondent (hispanic, black or white),
% sex: Sex of the respondent (male or female),
% ever_jailed: A boolean variable that equals 1 if the respondent was jailed between 1985 to 1994,
% composite_wealth_1985: A composite variable created by summing five wealth variables from the NLS data for the year 1985, 
% composite_wealth_1990: A composite variable created by summing five wealth variables from the NLS data for the year 1990, 
% composite_wealth_1996: A composite variable created by summing five wealth variables from the NLS data for the year 1996, 

% Goal: How did the wealth of individuals with a criminal record compare to those without in the years 1985, 1990, and 1996?

% Target hypothesis: In 1985, 1990, and 1996, individuals with a criminal record had significantly lower wealth compared to those without.
% """
% \end{minted}


% \begin{minted}[frame=lines,
% baselinestretch=1,
% bgcolor=Box3Color,
% fontsize=\footnotesize, breaklines, breaksymbolleft={}, breaksymbolright={}]{python}
% task =
% f"""\
% Load all datasets using python using provided paths.
% Paths: ['../DiscoveryBench/evolution_freshwater_fish/body-size-evolution-in-south-american-
% freshwater-fishes.csv'].

% Dataset name: body-size-evolution-in-south-american-freshwater-fishes.csv
% Dataset description: This dataset contains the drivers of speciation rates in South American freshwater fishes, employing an integrative approach that considers multiple biotic and abiotic factors.
% Brief description of columns:
% HYBAS_ID: Identifier for hydrological basins as defined by the HydroBASINS framework,
% long: The longitude for specific geographic locations used to map the presence/absence of species,
% lat: The latitude for specific geographic locations used to map the presence/absence of species,
% BAMM_speciation: This variable represents the average speciation rates of species within each sub-basin, as estimated by BAMM,
% BAMM_extinction: This variable represents the mean extinction rates for each tip of the phylogenetic tree, as estimated by the BAMM analysis,
% BAMM_NetDiv: This variable stands for the net diversification rates, which are calculated by subtracting the mean extinction rates from the mean 
% speciation rates for each tip in the phylogenetic tree,
% DR: Diversification Rate, a transformed measure of evolutionary distinctness to understand phylogenetic diversity. The lower the "ed" value, the higher the "DR" metric, indicating a species with less unique evolutionary history compared to others in the phylogeny,
% BEL_evol: Rates of Body Elongation evolution,
% MBL_evol: Rates of Maximum Body Length evolution,
% OGP_evol: Rates of Oral Gape Position Evolution,
% RES_evol: Rates of Relative Eye Size Evolution,
% RML_evol: Rate of Relative Maxillary Length evolution,
% bio1: This variable represents the annual mean temperature,
% bio12: This variable represents the annual mean precipitation,
% runoff: This variable is used to represent the quantity of water from precipitation that flows over the land's surface and does not get absorbed into the ground. It is being extracted from a geographic database (HydroAtlas)., aet: Actual evapotranspiration, which is the sum of evaporation and plant transpiration from the Earth's land and ocean surface to the atmosphere,
% Elevation: Mean elevation data indicating the height above sea level,
% sgr: Stream gradient is a measure of the steepness or incline of a stream or river. It can affect water flow and sediment transport, which in turn can influence the habitat conditions for freshwater species. Higher stream gradients usually correspond to faster-moving 
% water and can create different ecological conditions compared to lower gradients,
% soil_div: It measures the diversity of soil types or conditions within each sub-basin studied. Soil diversity is computed using a dataset that includes eight variables related to substrate types and soil conditions,
% area: The geographic area of a sub-basin, possibly in square meters, used as one of the predictors in the analysis., diversity: Species diversity, which in ecological studies typically refers to the number of species and their relative abundances in a given area, 

% Goal: What is the nature of the relationship between the rate of body elongation evolution and speciation rates?

% Target Hypothesis: The rate of body elongation evolution has no significant association with speciation rates.
% """
% \end{minted}

% % \begin{minted}[frame=lines,
% % baselinestretch=1,
% % bgcolor=Box3Color,
% % fontsize=\footnotesize, breaklines, breaksymbolleft={}, breaksymbolright={}]{python}
% % task =
% % requirements_engineering_for_ML_enabled_systems.8.0
% % f"""\
% % Load all datasets using python using provided paths.
% % Paths: ['../DiscoveryBench/requirements_engineering_for_ML_enabled_systems/
% % requirements_engineering_for_ML-enabled_systems.csv'].

% % Dataset name: requirements_engineering_for_ML-enabled_systems.csv
% % Dataset description: Survey responses detailing the roles, techniques, and documentation practices associated with requirements in ML-enabled system projects.
% % Brief description of columns:
% % ID: The unique identifier for each respondent.,
% % Status: The current status of the respondent,
% % Duration: The duration of the respondent's involvement,
% % D1_Undergraduation: Undergraduate (e.g., Computer Science, Statistics),
% % D1_Specialization: Specialization (e.g., Data Science specialization, Project Management specialization),
% % D1_Master: Master (e.g., M.Sc. in Computer Science, M.Sc. in Economics),
% % D1_Phd: Ph.D. (e.g., Ph.D. in Computer Science, Ph.D. in Mathematics),
% % D1_Courses: Professional ML Certifications/Courses (e.g., Google Professional ML Engineer Certification, Coursera/Udacity course on ML),
% % D1_Others: Other course specified by respondent,
% % D2_Country: Country in which the respondent is currently working,
% % D3_Company_Size: Size of the organization the respondent currently work for (1-10 employees, 11-50 employees ... more than 2000 employees),
% % D4_Role: Role that best describes the respondent's current activities within the company (Project Lead/ Project Manager, business Analyst, Requirements Engineer, Solution Architect, Data Scientist, Developer, Test Manager / Tester),
% % D4_Role_Others: Other role specified by respondent,
% % D5_Software_Experience: Years of experience in working with the development of software based products,
% % D6_ML_Experience: Years of Experience in developing ML-enabled systems,
% % D7_Total_ML_Projects: Number of ML-enabled system projects that the respondent participated in,
% % D8_ML_Production: Number of ML-enabled system projects that the respondent participated in that actually got deployed,
% % D9_ML_Project_Team_Size: The Team size of the ML-enabled system projects that the respondent participated in,
% % D10_ML_Management_Framework_None: Participant responded with None as the response for project management framework applied in the participated ML-enabled systems project,
% % D10_ML_Management_Framework_CRISP-DM: Participant responded with CRISP-DM as the response for project management framework applied in the participated ML-enabled systems project,
% % D10_ML_Management_Framework_Kanban: Participant responded with Kanban as the response for project management framework applied in the participated ML-enabled systems project,
% % D10_ML_Management_Framework_Lean: Participant responded with Lean as the response for project management framework applied in the participated ML-enabled systems project,
% % D10_ML_Management_Framework_RUP: Participant responded with RUP as the response for project management framework applied in the participated ML-enabled systems project,
% % D10_ML_Management_Framework_SAFe: Participant responded with SAFe as the response for project management framework applied in the participated ML-enabled systems project,
% % D10_ML_Management_Framework_Scrum: Participant responded with Scrum as the response for project management framework applied in the participated ML-enabled systems project,
% % D10_ML_Management_Framework_Others: Participant responded with a different framework as the response for project management framework applied in the participated ML-enabled systems project,
% % D10_ML_Management_Framework_Others_Free: Name of the other framework for project management framework applied in the participated ML-enabled systems project,
% % D11_Agile_Development: The agility of the development of the respondent in the ML-enabled systems projects that the respondent participated in,
% % D12_ML_Project_Context_Banking: Banking was the domain of application of the ML-enabled systems project that the respondent participated in,
% % D12_ML_Project_Context_Defense: Defense was the domain of application of the ML-enabled systems project that the respondent participated in,
% % D12_ML_Project_Context_Education: Education was the domain of application of the ML-enabled systems project that the respondent participated in,
% % D12_ML_Project_Context_Embedded: Embedded systems in Automotive or Avionics was the domain of application of the ML-enabled systems project that the respondent participated in,
% % D12_ML_Project_Context_Entertainment: Entertainment was the domain of application of the ML-enabled systems project that the respondent participated in,
% % D12_ML_Project_Context_Healthcare: Healthcare was the domain of application of the ML-enabled systems project that the respondent participated in,
% % D12_ML_Project_Context_Insurance: Insurance was the domain of application of the ML-enabled systems project that the respondent participated in,
% % D12_ML_Project_Context_Logistics: Logistics was the domain of application of the ML-enabled systems project that the respondent participated in,
% % D12_ML_Project_Context_Oil: Oil & Gas was the domain of application of the ML-enabled systems project that the respondent participated in,
% % D12_ML_Project_Context_Sales: Sales/E-commerce was the domain of application of the ML-enabled systems project that the respondent participated in,
% % D12_ML_Project_Context_Telecom: Telecommunication was the domain of application of the ML-enabled systems project that the respondent participated in,
% % D12_ML_Project_Context_Others: Respondent specified some other domain of application of the ML-enabled systems project that the respondent participated in,
% % D12_ML_Project_Context_Others_Free: Respondent's domain of application of the ML-enabled systems project that the respondent participated in,
% % D13_ML_Programming_Language_C: C language was in the list of general languages that composed the ML-enabled system projects (including eventually Non-ML related parts),
% % D13_ML_Programming_Language_Java: Java language was in the list of general languages that composed the ML-enabled system projects (including eventually Non-ML related parts),
% % D13_ML_Programming_Language_Javascript: Javascript language was in the list of general languages that composed the ML-enabled system projects (including eventually Non-ML related parts),
% % D13_ML_Programming_Language_Julia: Julia language was in the list of general languages that composed the ML-enabled system projects (including eventually Non-ML related parts),
% % D13_ML_Programming_Language_MatLab: MatLab language was in the list of general languages that composed the ML-enabled system projects (including eventually Non-ML related parts),
% % D13_ML_Programming_Language_Python: Python language was in the list of general languages that composed the ML-enabled system projects (including eventually Non-ML related parts),
% % D13_ML_Programming_Language_R: R language was in the list of general languages that composed the ML-enabled system projects (including eventually Non-ML related parts),
% % D13_ML_Programming_Language_Others: Other language was specified as the general language that composed the ML-enabled system projects (including eventually Non-ML related parts),
% % D13_ML_Programming_Language_Others_Free: Name of the other language that was specified as the general language that composed the ML-enabled system projects (including eventually Non-ML related parts),
% % D14_ML_Purpose_Prediction: Prediction was the main purpose of the ML-enabled system projects the respondent participated in,
% % D14_ML_Purpose_Prediction_Free: The typical purposes that were addressed using prediction in the project,
% % D14_ML_Purpose_Classification: Classification was the main purpose of the ML-enabled system projects the respondent participated in,
% % D14_ML_Purpose_Classification_Free: The typical purposes that were addressed using classification in the project,
% % D14_ML_Purpose_Association: Association was the main purpose of the ML-enabled system projects the respondent participated in,
% % D14_ML_Purpose_Association_Free: The typical purposes that were addressed using association in the project,
% % D14_ML_Purpose_Clustering: Clustering was the main purpose of the ML-enabled system projects the respondent participated in,
% % D14_ML_Purpose_Clustering_Free: The typical purposes that were addressed using clustering in the project,
% % D14_ML_Purpose_Others: ML-enabled system project had some other purpose,
% % D14_ML_Purpose_Others_Free: The other purposes that were addressed in the project,
% % D15_ML_Algorithms_Apriori: Apriori algorithm was employed in the ML-enabled system project that the respondent participated in,
% % D15_ML_Algorithms_Bayesian: Bayesian algorithm was employed in the ML-enabled system project that the respondent participated in,
% % D15_ML_Algorithms_DBSCAN: DBSCAN algorithm was employed in the ML-enabled system project that the respondent participated in,
% % D15_ML_Algorithms_Decision_Tree: Decision Tree algorithm was employed in the ML-enabled system project that the respondent participated in,
% % D15_ML_Algorithms_Ensembles: Ensemble  (e.g. Random Forests, XGBoost) Algorithm was employed in the ML-enabled system project that the respondent participated in,
% % D15_ML_Algorithms_Gaussian_Mixture: Gaussian Mixture was employed in the ML-enabled system project that the respondent participated in,
% % D15_ML_Algorithms_KMeans: KMeans algorithm was employed in the ML-enabled system project that the respondent participated in,
% % D15_ML_Algorithms_KNN: KNN was employed in the ML-enabled system project that the respondent participated in,
% % D15_ML_Algorithms_Linear_Regression: Linear Regression was employed in the ML-enabled system project that the respondent participated in,
% % D15_ML_Algorithms_Logistic_Regression: Logistic Regression was employed in the ML-enabled system project that the respondent participated in,
% % D15_ML_Algorithms_Naive_Bayes: Naive Bayes was employed in the ML-enabled system project that the respondent participated in,
% % D15_ML_Algorithms_Neural_Networks: Neural Networks were employed in the ML-enabled system project that the respondent participated in,
% % D15_ML_Algorithms_SVM: Support Vector Machines was employed in the ML-enabled system project that the respondent participated in,
% % D15_ML_Algorithms_Others: Some other algorithm was employed in the ML-enabled system project that the respondent participated in,
% % D15_ML_Algorithms_Others_Free: The name of the different algorithm that was employed in the ML-enabled system project that the respondent participated in,
% % Q1_ML_Life_Cycle_Importance_Problem_Understanding: The level of relevance of Problem Understanding and Requirements with regard to project success. One of the following: Not Relevant at All, Low Relevance, Neutral, High Relevance, Extremely Relevant, I don't know,
% % Q1_ML_Life_Cycle_Importance_Data_Collection: The level of relevance of Data Collection with regard to project success. One of the following: Not Relevant at All, Low Relevance, Neutral, High Relevance, Extremely Relevant, I don't know,
% % Q1_ML_Life_Cycle_Importance_Data_Pre-Processing: The level of relevance of Data Pre-Processing with regard to project success. One of the following: Not Relevant at All, Low Relevance, Neutral, High Relevance, Extremely Relevant, I don't know,
% % Q1_ML_Life_Cycle_Importance_Model_Creation: The level of relevance of Model Creation with regard to project success. One of the following: Not Relevant at All, Low Relevance, Neutral, High Relevance, Extremely Relevant, I don't know,
% % Q1_ML_Life_Cycle_Importance_Model_Evaluation: The level of relevance of Model Evaluation with regard to project success. One of the following: Not Relevant at All, Low Relevance, Neutral, High Relevance, Extremely Relevant, I don't know,
% % Q1_ML_Life_Cycle_Importance_Model_Deployment: The level of relevance of Model Deployment with regard to project success. One of the following: Not Relevant at All, Low Relevance, Neutral, High Relevance, Extremely Relevant, I don't know,
% % Q1_ML_Life_Cycle_Importance_Model_Monitoring: The level of relevance of Model Monitoring with regard to project success. One of the following: Not Relevant at All, Low Relevance, Neutral, High Relevance, Extremely Relevant, I don't know,
% % Q2_ML_Life_Cycle_Difficulty_Problem_Understanding: Difficulty level of Problem Understanding and Requirements stage in ML Life Cycle. One of the following: Very Easy, Easy, Neutral, Complex, Very Complex, I don't know,
% % Q2_ML_Life_Cycle_Difficulty_Data_Collection: Difficulty level of Data Collection stage in ML Life Cycle. One of the following: Very Easy, Easy, Neutral, Complex, Very Complex, I don't know,
% % Q2_ML_Life_Cycle_Difficulty_Data_Pre-Processing: Difficulty level of Data Pre-Processing stage in ML Life Cycle. One of the following: Very Easy, Easy, Neutral, Complex, Very Complex, I don't know,
% % Q2_ML_Life_Cycle_Difficulty_Model_Creation: Difficulty level of Model Creation stage in ML Life Cycle. One of the following: Very Easy, Easy, Neutral, Complex, Very Complex, I don't know,
% % Q2_ML_Life_Cycle_Difficulty_Model_Evaluation: Difficulty level of Model Evaluation stage in ML Life Cycle. One of the following: Very Easy, Easy, Neutral, Complex, Very Complex, I don't know,
% % Q2_ML_Life_Cycle_Difficulty_Model_Deployment: Difficulty level of Model Deployment stage in ML Life Cycle. One of the following: Very Easy, Easy, Neutral, Complex, Very Complex, I don't know,
% % Q2_ML_Life_Cycle_Difficulty_Model_Monitoring: Difficulty level of Model Monitoring stage in ML Life Cycle. One of the following: Very Easy, Easy, Neutral, Complex, Very Complex, I don't know,
% % Q3_ML_Life_Cycle_Effort_Problem_Understanding: The proportion of effort spent in the ML life cycle stage for Problem Understanding,
% % Q3_ML_Life_Cycle_Effort_Data_Collection: The proportion of effort spent in the ML life cycle stage for Data Collection,
% % Q3_ML_Life_Cycle_Effort_Data_Pre-Processing: The proportion of effort spent in the ML life cycle stage for Data Pre-Processing,
% % Q3_ML_Life_Cycle_Effort_Model_Creation: The proportion of effort spent in the ML life cycle stage for Model Creation,
% % Q3_ML_Life_Cycle_Effort_Model_Evaluation: The proportion of effort spent in the ML life cycle stage for Model Evaluation,
% % Q3_ML_Life_Cycle_Effort_Model_Deployment: The proportion of effort spent in the ML life cycle stage for Model Deployment,
% % Q3_ML_Life_Cycle_Effort_Model_Monitoring: The proportion of effort spent in the ML life cycle stage for Model Monitoring,
% % Q4_ML_Life_Cycle_Main_Problems_Problem_Understanding_Free_First: The first main problem faced in Problem Understanding phase in the ML life cycle stage,
% % Q4_ML_Life_Cycle_Main_Problems_Problem_Understanding_Free_Second: The second main problem faced in the Problem Understanding phase of the ML life cycle,
% % Q4_ML_Life_Cycle_Main_Problems_Problem_Understanding_Free_Third: The third main problem faced in the Problem Understanding phase of the ML life cycle,
% % Q4_ML_Life_Cycle_Main_Problems_Data_Collection_Free_First: The first main problem faced in the Data Collection phase of the ML life cycle,
% % Q4_ML_Life_Cycle_Main_Problems_Data_Collection_Free_Second: The second main problem faced in the Data Collection phase of the ML life cycle,
% % Q4_ML_Life_Cycle_Main_Problems_Data_Collection_Free_Third: The third main problem faced in the Data Collection phase of the ML life cycle,
% % Q4_ML_Life_Cycle_Main_Problems_Data_Pre-Processing_Free_First: The first main problem faced in the Data Pre-Processing phase of the ML life cycle,
% % Q4_ML_Life_Cycle_Main_Problems_Data_Pre-Processing_Free_Second: The second main problem faced in the Data Pre-Processing phase of the ML life cycle,
% % Q4_ML_Life_Cycle_Main_Problems_Data_Pre-Processing_Free_Third: The third main problem faced in the Data Pre-Processing phase of the ML life cycle,
% % Q4_ML_Life_Cycle_Main_Problems_Model_Creation_Free_First: The first main problem faced in the Model Creation phase of the ML life cycle,
% % Q4_ML_Life_Cycle_Main_Problems_Model_Creation_Free_Second: The second main problem faced in the Model Creation phase of the ML life cycle,
% % Q4_ML_Life_Cycle_Main_Problems_Model_Creation_Free_Third: The third main problem faced in the Model Creation phase of the ML life cycle,
% % Q4_ML_Life_Cycle_Main_Problems_Model_Evaluation_Free_First: The first main problem faced in the Model Evaluation phase of the ML life cycle,
% % Q4_ML_Life_Cycle_Main_Problems_Model_Evaluation_Free_Second: The second main problem faced in the Model Evaluation phase of the ML life cycle,
% % Q4_ML_Life_Cycle_Main_Problems_Model_Evaluation_Free_Third: The third main problem faced in the Model Evaluation phase of the ML life cycle,
% % Q4_ML_Life_Cycle_Main_Problems_Model_Deployment_Free_First: The first main problem faced in the Model Deployment phase of the ML life cycle,
% % Q4_ML_Life_Cycle_Main_Problems_Model_Deployment_Free_Second: The second main problem faced in the Model Deployment phase of the ML life cycle,
% % Q4_ML_Life_Cycle_Main_Problems_Model_Deployment_Free_Third: The third main problem faced in the Model Deployment phase of the ML life cycle,
% % Q4_ML_Life_Cycle_Main_Problems_Model_Monitoring_Free_First: The first main problem faced in the Model Monitoring phase of the ML life cycle,
% % Q4_ML_Life_Cycle_Main_Problems_Model_Monitoring_Free_Second: The second main problem faced in the Model Monitoring phase of the ML life cycle,
% % Q4_ML_Life_Cycle_Main_Problems_Model_Monitoring_Free_Third: The third main problem faced in the Model Monitoring phase of the ML life cycle,
% % Q4_ML_Life_Cycle_Main_Problems_Model_Other_Free_First: The first main problem faced in an unspecified phase of the ML life cycle,
% % Q4_ML_Life_Cycle_Main_Problems_Model_Other_Free_Second: The second main problem faced in an unspecified phase of the ML life cycle,
% % Q4_ML_Life_Cycle_Main_Problems_Model_Other_Free_Third: The third main problem faced in an unspecified phase of the ML life cycle,
% % Q5_ML_Life_Cycle_Main_Problems_Ranking_Free_First: The first main problem faced in ranking phase of the ML life cycle,
% % Q5_ML_Life_Cycle_Main_Problems_Ranking_Free_Second: The second main problem faced in ranking phase of the ML life cycle,
% % Q5_ML_Life_Cycle_Main_Problems_Ranking_Free_Third: The third main problem faced in ranking phase of the ML life cycle,
% % Q6_ML_Solution_Optimality: Degree to which the respondent believes ML solutions are optimal,
% % Q7_ML_Solution_Optimality_Extra_Effort: Extra effort required to achieve optimal ML solutions as perceived by the respondent,
% % Q8_ML_Addressing_Project_Lead: Degree to which ML aspects are addressed by the Project Lead in the respondent's organization,
% % Q8_ML_Addressing_Business_Analyst: Degree to which ML aspects are addressed by the Business Analyst in the respondent's organization,
% % Q8_ML_Addressing_Requirement_Engineer: Degree to which ML aspects are addressed by the Requirement Engineer in the respondent's organization,
% % Q8_ML_Addressing_Solution_Architect: Degree to which ML aspects are addressed by the Solution Architect in the respondent's organization,
% % Q8_ML_Addressing_Data_Scientist: Degree to which ML aspects are addressed by the Data Scientist in the respondent's organization,
% % Q8_ML_Addressing_Developer: Degree to which ML aspects are addressed by the Developer in the respondent's organization,
% % Q8_ML_Addressing_Tester: Degree to which ML aspects are addressed by the Tester in the respondent's organization,
% % Q8_ML_Addressing_Others: Degree to which ML aspects are addressed by other roles specified by the respondent,
% % Q8_ML_Addressing_Others_Free: Free text response for other roles addressing ML aspects specified by the respondent,
% % Q9_ML_Elicitation_Interviews: Degree to which interviews are used for ML requirements elicitation in the respondent's organization,
% % Q9_ML_Elicitation_Scenarios: Degree to which scenarios are used for ML requirements elicitation in the respondent's organization,
% % Q9_ML_Elicitation_Prototyping: Degree to which prototyping is used for ML requirements elicitation in the respondent's organization,
% % Q9_ML_Elicitation_Workshops_Meetings: Degree to which workshops and meetings are used for ML requirements elicitation in the respondent's organization,
% % Q9_ML_Elicitation_Observation: Degree to which observation is used for ML requirements elicitation in the respondent's organization,
% % Q9_ML_Elicitation_Others: Degree to which other methods are used for ML requirements elicitation specified by the respondent,
% % Q9_ML_Elicitation_Others_Free: Free text response for other methods of ML requirements elicitation specified by the respondent,
% % Q10_ML_Documentation_Not_Documented: Degree to which ML aspects are not documented in the respondent's organization,
% % Q10_ML_Documentation_Vision_Document: Degree to which vision documents are used for ML documentation in the respondent's organization,
% % Q10_ML_Documentation_Requirements_Lists: Degree to which requirements lists are used for ML documentation in the respondent's organization,
% % Q10_ML_Documentation_Goal_Models: Degree to which goal models are used for ML documentation in the respondent's organization,
% % Q10_ML_Documentation_Use_Case_Models: Degree to which use case models are used for ML documentation in the respondent's organization,
% % Q10_ML_Documentation_Prototypes: Degree to which prototypes are used for ML documentation in the respondent's organization,
% % Q10_ML_Documentation_User_Stories: Degree to which user stories are used for ML documentation in the respondent's organization,
% % Q10_ML_Documentation_BDD_Scenarios: Degree to which BDD scenarios are used for ML documentation in the respondent's organization,
% % Q10_ML_Documentation_MLCanvas: Degree to which ML Canvas is used for ML documentation in the respondent's organization,
% % Q10_ML_Documentation_Notebooks: Degree to which notebooks are used for ML documentation in the respondent's organization,
% % Q10_ML_Documentation_Data_Models: Degree to which data models are used for ML documentation in the respondent's organization,
% % Q10_ML_Documentation_Others: Degree to which other documentation methods are used for ML specified by the respondent,
% % Q10_ML_Documentation_Others_Free: Free text response for other documentation methods for ML specified by the respondent,
% % Q11_ML_NFRs_Not_Considered: Degree to which non-functional requirements (NFRs) are not considered in ML projects in the respondent's organization,
% % Q11_ML_NFRs_Data_Quality: Degree to which data quality is considered as a non-functional requirement in ML projects in the respondent's organization,
% % Q11_ML_NFRs_Model_Accountability: Degree to which model accountability is considered as a non-functional requirement in ML projects in the respondent's organization,
% % Q11_ML_NFRs_Model_Ethics_Fairness: Degree to which model ethics and fairness are considered as non-functional requirements in ML projects in the respondent's organization,
% % Q11_ML_NFRs_Model_Explainability: Degree to which model explainability is considered as a non-functional requirement in ML projects in the respondent's organization,
% % Q11_ML_NFRs_Model_Interactiveness: Degree to which model interactiveness is considered as a non-functional requirement in ML projects in the respondent's organization,
% % Q11_ML_NFRs_Model_Reliability: Degree to which model reliability is considered as a non-functional requirement in ML projects in the respondent's organization,
% % Q11_ML_NFRs_Model_Transparency: Degree to which model transparency is considered as a non-functional requirement in ML projects in the respondent's organization,
% % Q11_ML_NFRs_System_Compatibility: Degree to which system compatibility is considered as a non-functional requirement in ML projects in the respondent's organization,
% % Q11_ML_NFRs_System_Maintainability: Degree to which system maintainability is considered as a non-functional requirement in ML projects in the respondent's organization,
% % Q11_ML_NFRs_System_Performance: Degree to which system performance is considered as a non-functional requirement in ML projects in the respondent's organization,
% % Q11_ML_NFRs_System_Portability: Degree to which system portability is considered as a non-functional requirement in ML projects in the respondent's organization,
% % Q11_ML_NFRs_System_Privacy: Degree to which system privacy is considered as a non-functional requirement in ML projects in the respondent's organization,
% % Q11_ML_NFRs_System_Reliability: Degree to which system reliability is considered as a non-functional requirement in ML projects in the respondent's organization,
% % Q11_ML_NFRs_System_Safety: Degree to which system safety is considered as a non-functional requirement in ML projects in the respondent's organization,
% % Q11_ML_NFRs_System_Security: Degree to which system security is considered as a non-functional requirement in ML projects in the respondent's organization,
% % Q11_ML_NFRs_System_Usability: Degree to which system usability is considered as a non-functional requirement in ML projects in the respondent's organization,
% % Q11_ML_NFRs_Others: Degree to which other non-functional requirements are considered in ML projects specified by the respondent,
% % Q11_ML_NFRs_Others_Free: Free text response for other non-functional requirements considered in ML projects specified by the respondent,
% % Q12_ML_Most_Difficult_Activity_Customer_Expectations: Difficulty in managing customer expectations in ML projects as perceived by the respondent,
% % Q12_ML_Most_Difficult_Activity_Eliciting_Analyzing: Difficulty in eliciting and analyzing requirements in ML projects as perceived by the respondent,
% % Q12_ML_Most_Difficult_Activity_Aligning_Requirements_Data: Difficulty in aligning requirements with data in ML projects as perceived by the respondent,
% % Q12_ML_Most_Difficult_Activity_Conflicts: Difficulty in resolving conflicts in ML projects as perceived by the respondent,
% % Q12_ML_Most_Difficult_Activity_New_Quality_Attributes: Difficulty in dealing with new quality attributes in ML projects as perceived by the respondent,
% % Q12_ML_Most_Difficult_Activity_Documentation: Difficulty in documenting ML projects as perceived by the respondent,
% % Q12_ML_Most_Difficult_Activity_Selecting_Metrics: Difficulty in selecting appropriate metrics in ML projects as perceived by the respondent,
% % Q12_ML_Most_Difficult_Activity_Verification: Difficulty in verifying ML projects as perceived by the respondent,
% % Q12_ML_Most_Difficult_Activity_Changing_Requirements: Difficulty in managing changing requirements in ML projects as perceived by the respondent,
% % Q12_ML_Most_Difficult_Activity_Others: Difficulty in other activities in ML projects specified by the respondent,
% % Q12_ML_Most_Difficult_Activity_Others_Free: Free text response for other difficult activities in ML projects specified by the respondent,
% % Q13_Model_Deploy_Approach_Embedded_Model: Degree to which the embedded model deployment approach is used in the respondent's organization,
% % Q13_Model_Deploy_Approach_Service: Degree to which the service model deployment approach is used in the respondent's organization,
% % Q13_Model_Deploy_Approach_PaaS: Degree to which the Platform as a Service (PaaS) deployment approach is used in the respondent's organization,
% % Q13_Model_Deploy_Approach_Others: Degree to which other deployment approaches are used in the respondent's organization,
% % Q13_Model_Deploy_Approach_Others_Free: Free text response for other deployment approaches used in the respondent's organization,
% % Q14_Model_Deploy_Pipeline_Yes: Yes response indicating if a deployment pipeline is used in the respondent's organization,
% % Q14_Model_Deploy_Pipeline_Yes_Free: Free text response if a deployment pipeline is used in the respondent's organization,
% % Q14_Model_Deploy_Pipeline_No: No response indicating if a deployment pipeline is not used in the respondent's organization,
% % Q15_Model_Deploy_Production_Monitoring: Degree to which production monitoring is conducted for deployed models in the respondent's organization,
% % Q16_Model_Monitor_Aspects_Input_And_Output: Importance of monitoring inputs and outputs of models in the respondent's organization,
% % Q16_Model_Monitor_Aspects_Interpretability_Output: Importance of monitoring the interpretability of model outputs in the respondent's organization,
% % Q16_Model_Monitor_Aspects_Output_And_Decisions: Importance of monitoring outputs and decisions of models in the respondent's organization,
% % Q16_Model_Monitor_Aspects_Fairness: Importance of monitoring fairness of models in the respondent's organization,
% % Q16_Model_Monitor_Aspects_Others: Importance of monitoring other aspects of models specified by the respondent,
% % Q16_Model_Monitor_Aspects_Others_Free: Free text response for other aspects of model monitoring specified by the respondent,
% % Q17_Automated_Machine_Learning_Tools_Yes_No: Yes or No response indicating if the respondent uses automated machine learning tools,
% % Q17_Automated_Machine_Learning_Tools_Yes_Free: Free text response if the respondent uses automated machine learning tools,
% % Origin: Origin of the respondent,

% % Goal: Which two documentation formats are the least used for requirements in ML-enabled system projects, with 10.13% (95% CI [9.926, 10.333]) and 4.366% (95% CI [4.231, 4.501]) of respondents indicating so, respectively, after bootstrapping for statistical significance?

% % Hypothesis:
% % """
% % \end{minted}

% \section{\synth{} Task Examples}
% \label{app:synth-examples}

% \begin{minted}[frame=lines,
% baselinestretch=1,
% bgcolor=Box3Color,
% fontsize=\footnotesize, breaklines, breaksymbolleft={}, breaksymbolright={}]{python}
% task =
% f"""\
% Load all datasets using python using provided paths.
% Paths: ['../DiscoveryBenchSynth/benchmark_split/test/ancient-languages_0_2/data.csv'].

% Dataset name: data.csv
% Dataset description: Related to hieroglyphs, Sanskrit, Latin, and deciphering old scripts.
% Brief description of columns:
% digital_resources_availability: Indicates te amount of digital resources available, rated from 1 (least) to 5 (most),
% language_family_rank: Ranking of language families based on the commonality of study, where a smaller number indicates a more commonly studied language family,
% geopolitical_tension: A measure of current geopolitical tensions in the region from where the script originated, on a scale from 1 to 10,
% deciphering_tools_count: Number of tools available specifically for deciphering scripts of each ancient language family,
% year_of_discovery: The year in which the script was discovered,
% use_of_ml: Indicates if machine learning methods were used in the deciphering process, 
% artifact_density: Count of artifacts found per square kilometer in the immediate area of the script discovery,
% script_complexity: Complexity rating of the ancient script on a scale from 1 to 10, with 10 being the most complex,
% script_complexity_level: The complexity level of the script to be translated, rated from 1 (least complex) to 5 (most complex),
% is_currently_used: Indicates whether the script is still in use today,
% material_complexity: The complexity of the material found in the scripts based on linguistic and symbolic features,
% preservation_quality: Assessment of the script's condition on a scale from 1 (poor) to 10 (excellent),
% script_distinctiveness: A measure of how unique the script is compared to others, on a scale from 1 to 10,
% script_characteristic_importance: Importance of the script's characteristics on a scale from 1 (low) to 5 (high), e.g, uniqueness or historical significance,
% technological_support: Indicates if advanced technological support (machine learning tools) is used (1) or not (0),
% high_accuracy_requirement_percentage: The percentage of translation projects that need very high accuracy,
% recognized_in_cultural_practices: Whether elements of the ancient language are recognized and used in current cultural or societal practices,
% accuracy_percentage: The percentage of accuracy in the translations achieved by the method, 
% deciphering_method: The method used for deciphering scripts: traditional or machine learning, 
% included_in_academic_curriculum: Indicates whether the language is part of the academic curriculum in significant number of higher education institutions (more than 50% globally), 
% government_endorsement: Whether the language receives official recognition and support from the government,
% media_representation: The number of appearances or features in mainstream media outlets over the past year,
% estimated_origin_year: Estimated year when the language first originated, 
% specialized_researchers_count: Number of researchers that specialize in each ancient language family,
% sediment_accumulation_rate: Average rate of sediment accumulation at the discovery site, in millimeters per year,
% has_modern_derivative: Indicates if the ancient language has a modern derivative that is still in use,
% international_collaboration: Indicates if there is any international collaboration in the translation project,
% academic_conferences_count: Number of academic conferences held annually dedicated to the language,
% vocabulary_size: The estimated number of unique words used in the ancient language, 
% expert_translators_available: The number of expert translators available for the project, 
% language_family: The family to which the ancient language belongs, 
% educational_program_availability: Indicates if structured educational programs are available for the language,
% has_dedicated_funding: Indicates whether there is dedicated funding available for the translation efforts of each language family,
% collaborative_effort: Indicates whether the deciphering was a collaborative effort between multiple experts or not,
% recent_publications: Indicates if there has been a recent publication in the field of ancient languages within the last month,
% archaeology_technology_index: A quantitative index of the technological advancement of archaeological tools used in the excavation, on a scale from 1 to 100, 
% foundational_for_modern_languages: The count of modern languages that are directly descended from or heavily influenced by the ancient language,
% digitized_artifacts_count: The number of digitized artifacts available for each language family,
% international_experts_involved: Whether international experts were involved in the translation process, true for involved, false for not involved,
% active_researchers_count: Number of active researchers (academic and amateurs) dedicated to studying the language,
% percentage_of_deciphered_texts: The percentage of known texts that have been successfully deciphered,

% Goal: Is there a relationship between the involvement of international experts in analysing ancient languages and the level of geopolitical tension in various regions?

% Target hypothesis: The involvement of international experts is more likely when there is high geopolitical tension in the region, specifically in the Middle East.
% """
% \end{minted}


% \begin{minted}[frame=lines,
% baselinestretch=1,
% bgcolor=Box3Color,
% fontsize=\footnotesize, breaklines, breaksymbolleft={}, breaksymbolright={}]{python}
% task =
% f"""\
% Load all datasets using python using providd paths.
% Paths: ['../DiscoveryBenchSynth/benchmark_split/test/astronomy_0_3/data.csv'].

% Dataset name: data.csv
% Dataset description: Related to celestial bodies, galaxies, and space phenomena.
% Brief description of columns:
% distance_to_nearest_star_forming_region: Distance from the galactic center to the nearest star-forming region measured in light-years,
% star_density: Number of stars per unit volume in the galaxy,
% passive_star_nuclei: Count of galaxies in a cluster with passive or dormant galactic nuclei,
% blue_signature_intensity: Measure of blue light intensity from the galaxy,
% gravitational_pull: The combined gravitational force exerted by nearby celestial bodies, measured in multiples of Earth's gravity,
% star_formation_rate: Rate of star formation in the galaxy measured in solar masses per year,
% gravity_rule_index: A score indicating the strength of gravitational pull within the galaxy,
% stellar_density: Number of stars per cubic light year in the galaxy,
% dominant_star_type: Most common type of star within the galaxy,
% orbital_companions: Number of smaller celestial bodies or satellites orbiting the galaxy,
% average_star_surface_temperature: Average temperature on the surface of stars in the galaxy measured in Kelvin,
% galaxy_mass: The total mass of the galaxy measured in solar masses,
% cosmic_ray_exposure: Measure of the intensity of cosmic ray exposure in particles per square meter,
% number_of_satellite_galaxies: The count of smaller galaxies gravitationally bound to the main galaxy,
% central_black_hole_presence: Indicates whether there is a supermassive black hole at the center of the galaxy,
% high_radiation_levels: Binary indicator of whether the galaxy is exposed to high levels of cosmic radiation,
% recent_supernova_events: Number of supernova events in the galaxy's quadrant within the past million years,
% average_surface_temperature_of_stars: Average surface temperature of the dominant stars within the galaxy, in Kelvin,
% supernova_rate: Number of supernovae observed per year in the galaxy,
% cosmic_radiation_level: Average level of cosmic radiation in the galaxy, measured in microsieverts per hour (µSv/h),
% average_orbital_velocity: The average velocity at which celestial bodies orbit the center of the galaxy, measured in kilometers per second,
% metallic_content_percentage: Average percentage of metals found in stars within the galaxy,
% has_active_galactic_nucleus: Indicates whether the galaxy has an active galactic nucleus,
% has_binary_star_system: Binary status indicating if the galaxy contains a binary star system,
% galaxy_type: Classification of the galaxy based on its structure,
% colder_cosmic_temperature: Temperature of the cosmic environment of the galaxy measured in Kelvin,
% intense_stellar_formation: Binary indicator of whether the galaxy's cluster has a high rate of star formation,
% galaxy_ID: Unique identifier for each galaxy,
% ultraviolet_emission: Indicator whether the galaxy has significant ultraviolet emission, binary true/false,
% binary_star_percentage: The percentage of stars within a galaxy that are part of a binary system,
% presence_exotic_matter: Binary status indicating if exotic matter is significantly present within the galaxy,
% cosmic_radiation_percentile: Relative percentile rank of a galaxy based on detected cosmic radiation levels,
% interstellar_magnetic_field_strength: Strength of the galaxy's interstellar magnetic field, measured in gauss,
% dark_matter_percentage: The percentage of a galaxy's mass that is made up of dark matter,
% velocity_dispersion_of_stars: Statistical spread in velocities of stars within the galaxy, in km/s,
% binary_star_systems: Number of binary star systems observed in the galaxy,
% proportion_of_red_stars: Proportion of red stars compared to total stars in the galaxy,
% active_star_nuclei: Count of galaxies in a cluster with active galactic nuclei,
% nova_light_emission_intensity: Mean intensity of light emissions from nova events near the galaxy measured in lumens,
% luminosity_variability: Fluctuation degree of galaxy's luminosity over a 10-year period, measured as the standard deviation of luminosity,
% central_black_hole_mass: Mass of the central black hole in the galaxy, measured in solar masses,
% galaxy_age: Estimated age of the galaxy in billions of years,
% variable_stellar_types_ratio: The ratio of variable star types to total star count within the galaxy,
% magnetic_field_strength: Measure of the galaxy's average magnetic field strength in gauss,
% dark_matter_anomaly_gravitational_effect: Measured gravitational effect of the nearest dark matter anomaly on the galaxy, on a scale from weak to strong,
% star_forming_region_proximity: Closeness of the galaxy to nearest star-forming regions, measured in light years,
% dominant_stellar_age: Average age of the dominant stellar population in the galaxy, measured in millions of years,
% binary_star_systems_percentage: Percentage of star systems in the galaxy that are binary or multiple star systems,
% dark_matter_density_classification: Binary classification of dark matter density in the region,
% is_in_orion_cluster: Binary status indicating if the galaxy is located in the Orion cluster,
% magnetic_field_presence: Indicator of the presence of significant magnetic fields in the galaxy,
% distance_to_galactic_core: Distance of the galaxy from the nearest galactic core in light years,
% galaxy_spectral_redshift: Measure of the galaxy's redshift indicating the speed at which it is moving away,
% number_of_nebulae: Count of the nebulae within the galaxy,
% interstellar_gas_mass: Total mass of interstellar gas in the galaxy measured in solar masses,
% distance_to_galactic_center: The distance of the galaxy from the galactic center in light years,
% cosmic_structures: Number of larger cosmic structures within the vicinity of the galaxy,
% proximity_to_black_hole: Distance of the galaxy from the nearest black hole in light years, 

% Goal: Is there a relationship between the intensity of light emissions from nova events near a galaxy and the percentage of metallic content in stars, that could be indicative of the gravitational effect of a nearby dark matter anomaly?

% Target Hypothesis: The intensity of light emissions from nova events near the galaxy and the percentage of metallic content in stars are used to determine the gravitational effect of the nearest dark matter anomaly on the galaxy, from weak to strong.
% """
% \end{minted}

% \begin{minted}[frame=lines,
% baselinestretch=1,
% bgcolor=Box3Color,
% fontsize=\footnotesize, breaklines, breaksymbolleft={}, breaksymbolright={}]{python}
% task =
% f"""\
% Load all datasets using python using provided paths.
% Paths: ['../DiscoveryBenchSynth/benchmark_split/test/futuristic-technology_1_1/data.csv'].

% Dataset name: data.csv
% Dataset description: Related to cutting-edge inventions, future gadgets, and sci-fi concepts.
% Brief description of columns: household_income: Categorical income levels of the household,
% age_of_household_head: Age of the head of the household,
% flagship_tech_company_presence: Indicates if there is at least one flagship technology company (e.g., major corporations in innovation) in the area,
% tech_workshops_frequency: Number of technology-based workshops conducted in the region per year,
% robot_programming_skill_level: Rated skill level of robot programming in the household from beginner (1), intermediate (2), to advanced (3),
% year: The year of observation,
% digital_infrastructure_index: An index from 0 to 10 measuring the quality and accessibility of digital infrastructure such as internet speed, public Wi-Fi access points, and digital education tools,
% urban_area: Whether the household is in an urban area,
% public_interest: A rating from 0 to 100 reflecting public interest in adopting futuristic technologies in the area,
% tech_patent_applications_per_1000: Number of technological patent applications filed per 1000 residents in the area annually,
% investment_in_tech_sector_percentage: Percentage of the family's total investments that are allocated to technology sectors,
% government_investment: Amount of funds invested by the government in futuristic technology per year in millions,
% number_of_robots: Number of robot assistants used in the household,
% households_with_internet_percent: Percentage of households in the area that have internet access,
% smart_home_ownership: Indicates if the household has smart home systems installed,
% public_interest_in_tech: Binary indicator of public interest in technology, where 1 represents high interest and 0 represents low interest, based on surveys and media analysis,
% stem_graduate_percentage: Percentage of the total population that graduates each year with a degree in science, technology, engineering, or mathematics,
% VR_training_facility_present: Indicates whether the household has a virtual reality setup for technical training,
% tech_openness_score: Score rating the household's openness to technology on a scale from 1 to 10,
% regulatory_change_rate: The count of significant technology-related regulatory updates made per year.,
% annual_tech_expos: Counts the number of technology exhibitions held annually in the area,
% decision_making_robots: Whether robots used can perform tasks requiring human decision-making,

% Goal: Is there a relationship between the age of the household head and a combination of the technological openness score, smart home ownership, and robot programming skill level? If so, what is this relationship and can the age be predicted based on these factors?

% Target Hypothesis: The age of the head of the household can be predicted by subtracting two times the technological openness score, five times the ownership of smart home systems, and ten times the skill level of robot programming from 100. Higher values of these factors result in a lower predicted age of the household head.
% """
% \end{minted}


% % \tushar{Good to keep a track of the size of supplementary - both the size and pages. Pages would not matter much if they are at the end.}
% % \ashish{add ellipsis where possible in the context. small is big! people should get an idea looking at diversity of data}

% % \bodhi{we did not add the details of croissant and huggingface download as we were waiting for review feedback - as our data structure is changing}

% % \bodhi{Longterm data governance and repo sustainability only. No need to keep on living repo.}


% % \begin{small}
% % \begin{verbatim}
% % Usage: discovery_agent.py [OPTIONS] QUERY

% % Options:
% %   --agent_type [coder|react]    Agent type to use for discovery, default is
% %                                 coder
% %   --model_name TEXT             Model name, default is gpt-4o, available
% %                                 models are [gpt-4-turbo|llama-3-70b-chat|claud
% %                                 e-3-opus|gemini-pro]. Exhaustive list can be
% %                                 found in config/model_config.json
% %   --api_config TEXT             API config file, default is
% %                                 config/api_config.json
% %   --log_file TEXT               Log file
% %   --metadata_path TEXT          Metadata file path  [required]
% %   --metadata_type [real|synth]  Metadata type  [required]
% %   --add_domain_knowledge        Add domain knowledge to query
% %   --add_workflow_tags           Add Workflow Tags to query
% %   --help                        Show this message and exit.
% % \end{verbatim}
% % \end{small}

% \section{Example tasks by domain}

% % \subsection{WIP section for finalizing prompt details}
% % \begin{enumerate}
% %     \item Sociology: ex1: nlsraw-6-0, ex2: nls-ses-10-1, ex3: nls-incarceration-3-3

    
    
% %     \item Biology: evolution-3-0, introduction-pathways-4-2, introduction-pathways-0-2, 
% %     \item Economics: ex1: worldbankeducation-gdp-2-0, ex2: worldbank-education-gdp-indicators-4-0, ex3: immigration-1-1
% %     \item Engineering: ex1: ML-enabled-systems-14-0
% %     \item Meta-science: meta-raw-0-3,  meta-raw-19-0, meta-2-0
% %     \item Humanities (with domain knowledge): 
% % \end{enumerate}

% % \paragraph{
% Note: Some of the dataset details are long. We trimmed them by adding an ellipsis for easy viewing. Examples are denoted by: folder\_name $\rightarrow$ metadata\_id $\rightarrow$ qid.
% % }

% \subsection{Sociology}

% Example 1: nls\_raw $\rightarrow$ 6 $\rightarrow$ 0

% \begin{minted}[frame=lines,
% baselinestretch=1,
% bgcolor=Box3Color,
% fontsize=\footnotesize, breaklines, breaksymbolleft={}, breaksymbolright={}]{python}
% task =
% f"""\

% Load all datasets using python using provided paths.
% Paths: ['../DiscoveryBench/nls_raw/nls_raw.csv']

% Dataset name: nls_raw.csv
% Dataset description: The dataset contains information from National Longitudinal Survey of Youth (NLSY79). It includes information about the Demographics, Family Background, Education, Health, Residential, Financial & Criminal Records of the participants.
% Brief description of columns:
% ID# (range 1-12686) 1979: Unique Identifier of the respondent,
% Sample ID, 1979 (interview): Sample Identification Code,
% Age of respondent, 1979: Age of respondent in 1979,
% Age of respondent at interview date, 1981: Age of respondent in 1981,
% Age of respondent at interview date, 1989: Age of respondent in 1989,
% Occupation of adult male in household at age 14, 1979: Occupation of the adult male present in the household of the respondent at age 14 in 1979. Variable records the occupation of the father figure of the repondent, values include FARMER AND FARM MANAGERS, PROFESSIONAL,TECHNICAL AND KINDRED etc,
% Highest grade completed by respondent's mother, 1979: Highest grade or year of regular school that respondent's mother ever completed till 1979,
% Highest grade completed by respondent's father, 1979: Highest grade or year of regular school that respondent's father ever completed till 1979,
% Highest grade completed, 1979: Highest grade or year of regular school that respondent have completed and got credit for till 1979,
% Racial/ethnic cohort, 1979: Respondent's racial/ethnic cohort, contains one of three values 1:BLACK, 2:HISPANIC, 3:NON-BLACK NON-HISPANIC,
% Sex of respondent, 1979: Sex of the respondent, 1:MALE or 2:FEMALE,
% Family size, 1979: Family size of the respondent in 1979,
% Ever convicted of an illegal act in adult court before 1980: Boolean variable that indicates if the respondent was convicted of an illegal act in adult court other than minor traffic violations before 1980,
% Ever been sentenced in any correctional institution before 1980: Boolean variable that indicated if the respondent was sentenced to spend time in a corrections institute, like a jail, prison, or a youth institution like a training school or reform school or not before 1980,
% Height of respondent, 1981: Height of the respondent in inches in 1981,
% Height of respondent, 1985: Height of the respondent in inches in 1985,
% Weight of respondent, 1981: Weight of the respondent in kilograms in 1981,
% Weight of respondent, 1989: Weight of the respondent in kilograms in 1989,
% Weight of respondent, 1992: Weight of the respondent in kilograms in 1992,
% Rank in class last year attended at this school, 1981: Respondent's rank in the class that he attended in school last year (in 1980) (variable recorded in 1981),
% Number of students in class last year attended at this school, 1981: Number of students in the respondent's class for the last year attended this school,
% ASVAB - Arithmetic Reasoning Z Score (rounded), 1981: This variable represents the standardized scores of respondents on the Arithmetic Reasoning section of the ASVAB test. It provides a way to compare individuals' performance on this specific aspect of the test within a standardized framework.,
% ASVAB - Word Knowledge Z Score (rounded), 1981: This variable represents the standardized scores of respondents on the Word Knowledge section of the ASVAB test, allowing for comparison of individuals' performance on this specific aspect of the test within a standardized framework.,
% ASVAB - Paragraph Comprehension Z Score (rounded), 1981: This variable represents the standardized scores of respondents on the Paragraph Comprehension section of the ASVAB test, allowing for comparison of individuals' performance on this specific aspect of the test within a standardized framework.,

% ASVAB - Mathematics Knowledge Z Score (rounded), 1981: This variable represents the standardized scores of respondents on the Mathematics Knowledge section of the ASVAB test, facilitating comparison of individuals' performance on this specific aspect of the test within a standardized framework.,
% Type of residence respondent is living in, 1981: Type of residence respondent is living in the 1981, contains one of these values 1:ABOARD SHIP, BARRACKS,    2:BACHELOR, OFFICER QUARTERS,    3:DORM, FRATERNITY, SORORITY,    4:HOSPITAL,    5:JAIL,    6:OTHER TEMPORARY QUARTERS,    11:OWN DWELLING UNIT,    12:ON-BASE MIL FAM HOUSING,    13:OFF-BASE MIL FAM HOUSING,    14:ORPHANAGE,    15:RELIGIOUS INSTITUTION,    16:OTHER INDIVIDUAL QUARTERS,    17:PARENTAL,    18:HHI CONDUCTED WITH PARENT,    19:R IN PARENTAL HOUSEHOLD,
% Type of residence respondent is living in, 1982: Type of residence respondent is living in the 1982, contains one of these values 1:ABOARD SHIP, BARRACKS,    2:BACHELOR, OFFICER QUARTERS,    3:DORM, FRATERNITY, SORORITY,    4:HOSPITAL,    5:JAIL,    6:OTHER TEMPORARY QUARTERS,    11:OWN DWELLING UNIT,    12:ON-BASE MIL FAM HOUSING,    13:OFF-BASE MIL FAM HOUSING,    14:ORPHANAGE,    15:RELIGIOUS INSTITUTION,    16:OTHER INDIVIDUAL QUARTERS,    17:PARENTAL,    18:HHI CONDUCTED WITH PARENT,    19:R IN PARENTAL HOUSEHOLD,
% ...
% ...
% ...
% Type of residence respondent is living in, 1993: Type of residence respondent is living in the 1993, contains one of these values 1:ABOARD SHIP, BARRACKS,    2:BACHELOR, OFFICER QUARTERS,    3:DORM, FRATERNITY, SORORITY,    4:HOSPITAL,    5:JAIL,    6:OTHER TEMPORARY QUARTERS,    11:OWN DWELLING UNIT,    12:ON-BASE MIL FAM HOUSING,    13:OFF-BASE MIL FAM HOUSING,    14:ORPHANAGE,    15:RELIGIOUS INSTITUTION,    16:OTHER INDIVIDUAL QUARTERS,    17:PARENTAL,    18:HHI CONDUCTED WITH PARENT,    19:R IN PARENTAL HOUSEHOLD,
% Type of residence respondent is living in, 1994: Type of residence respondent is living in the 1994, contains one of these values 1:ABOARD SHIP, BARRACKS,    2:BACHELOR, OFFICER QUARTERS,    3:DORM, FRATERNITY, SORORITY,    4:HOSPITAL,    5:JAIL,    6:OTHER TEMPORARY QUARTERS,    11:OWN DWELLING UNIT,    12:ON-BASE MIL FAM HOUSING,    13:OFF-BASE MIL FAM HOUSING,    14:ORPHANAGE,    15:RELIGIOUS INSTITUTION,    16:OTHER INDIVIDUAL QUARTERS,    17:PARENTAL,    18:HHI CONDUCTED WITH PARENT,    19:R IN PARENTAL HOUSEHOLD,
% Type of residence respondent is living in, 1996: Type of residence respondent is living in the 1996, contains one of these values 1:ABOARD SHIP, BARRACKS,    2:BACHELOR, OFFICER QUARTERS,    3:DORM, FRATERNITY, SORORITY,    4:HOSPITAL,    5:JAIL,    6:OTHER TEMPORARY QUARTERS,    11:OWN DWELLING UNIT,    12:ON-BASE MIL FAM HOUSING,    13:OFF-BASE MIL FAM HOUSING,    14:ORPHANAGE,    15:RELIGIOUS INSTITUTION,    16:OTHER INDIVIDUAL QUARTERS,    17:PARENTAL,    18:HHI CONDUCTED WITH PARENT,    19:R IN PARENTAL HOUSEHOLD,
% Family net wealth, 1985: Total Net Wealth for Family. Created by summing all asset values and subtracting all debts for the year 1985,
% Family net wealth, 1990: Total Net Wealth for Family. Created by summing all asset values and subtracting all debts for the year 1990,
% Family net wealth, 1996 (key data point): Total Net Wealth for Family. Created by summing all asset values and subtracting all debts for the year 1996,
% Market value of residential property respondent/spouse own, 1985: Market value of residential property that respondent/spouse owned in 1985,
% Market value of residential property respondent/spouse own, 1990: Market value of residential property that respondent/spouse owned in 1990,
% Market value of residential property respondent/spouse own, 1996: Market value of residential property that respondent/spouse owned in 1996,
% Total market value of farm, business, and other property, 1985: Total market value of all of the real estate, assets in the business(es), farm operation(s) in 1985,
% Total market value of farm, business, and other property, 1990: Total market value of all of the real estate, assets in the business(es), farm operation(s) in 1990,
% Total market value of farm, business, and other property, 1996: Total market value of all of the real estate, assets in the business(es), farm operation(s) in 1996,
% Market Value of vehicles respondent/spouse own, 1985: Total market value of all vehicles including automobiles that respondent/spouse owned in 1985,
% Market Value of vehicles respondent/spouse own, 1990: Total market value of all vehicles including automobiles that respondent/spouse owned in 1990,
% Market Value of vehicles respondent/spouse own, 96: Total market value of all vehicles including automobiles that respondent/spouse owned in 1996,
% Total market value of items over $500, 1985: Total market value of all the other assets of the respondent that were worth more than $500 in 1985,
% Total market value of items over $500, 1990: Total market value of all the other assets of the respondent that were worth more than $500 in 1990,
% Total market value of items over $500, 1996: Total market value of all the other assets of the respondent that were worth more than $500 in 1996,
% Total net family income, previous calendar year, 1979: Total net family income for the previous calendar year (1978) (recorded in 1979),
% Total net family income, previous calendar year, 1985: Total net family income for the previous calendar year (1984) (recorded in 1985),
% Total net family income, previous calendar year, 1989: Total net family income for the previous calendar year (1989) (recorded in 1989),
% Was more money put into or taken out of R/spouse savings since last interview, 1989: Categorical variable indicating if was more money was put into or taken out of respondent/spouse savings since last interview in 1989.
% It contains four values 1:PUT MORE MONEY IN, 2:TOOK MORE MONEY OUT, 3:NO CHANGE, 4:NO SAVINGS,
% Net amount respondent/spouse put into savings since last interview, 1989: Net amount of money that respondent/spouse put into their savings since last interview in 1989,
% Net amount respondent/spouse took out of savings since last interview, 1989: Net amount of money that respondent/spouse took out of savings since last interview in 1989,

% Goal: How is the advantage in BA degree completion rates for Black students related to Socioeconomic status levels?

%         \end{minted}


% Example 2: nls\_ses $\rightarrow$ 10 $\rightarrow$ 1

% \begin{minted}[frame=lines,
% baselinestretch=1,
% bgcolor=Box3Color,
% fontsize=\footnotesize, breaklines, breaksymbolleft={}, breaksymbolright={}]{python}
% task =
% f"""\

% Load all datasets using python using provided paths.
% Paths: ['../DiscoveryBench/nls_ses_processed/nls_ses_processed.csv']

% Dataset name: nls_ses_processed.csv
% Dataset description: This dataset contains social background factors (race, gender, and socioeconomic status) and academic resources at the time of secondary school graduation (standardized test scores, class rank, and curriculum) for the participants of the NLS.
% Brief description of columns:
% CASE ID: Unique ID of each respondent,
% SAMPLE_RACE: Race of the respondent (Hispanic, Black or White),
% SAMPLE_SEX: Sex of the respondent (Male or Female),
% FAMILY SIZE OF SAMPLE: Family size of the respondent,
% ABILITY: COMPOSITE OF ASVAB SCORE: Composite variable created by summing following variables:

% ASVAB - Arithmetic Reasoning Z Score (rounded), 1981
% ASVAB - Word Knowledge Z Score (rounded), 1981
% ASVAB - Paragraph Comprehension Z Score (rounded), 1981
% ASVAB - Mathematics Knowledge Z Score (rounded), 1981,
% BA DEGREE COMPLETED: Boolean variable that equals 1 if the BA Degree was completed by the respondent else 0,
% PERCENTILE IN CLASS: Respondent's percentile in the class that he attended in school last year,
% SES: Socioeconomic Status of the respondent,

% Goal: How does the effect of race on BA Degree completion change when both SES and academic characteristics are considered as compared to when only SES is considered?
% """
% \end{minted}

% Example 3: nls\_incarceration $\rightarrow$ 10 $\rightarrow$ 1

% \begin{minted}[frame=lines,
% baselinestretch=1,
% bgcolor=Box3Color,
% fontsize=\footnotesize, breaklines, breaksymbolleft={}, breaksymbolright={}]{python}
% task =
% f"""\
% Load all datasets using python using provided paths.
% Paths: ['../DiscoveryBench/nls_incarceration_processed/nls_incarceration_processed.csv']

% Dataset name: nls_incarc_processed.csv
% Dataset description: This dataset was created from the National Longitudinal Study of Youth's 1979 cohort to about the race, wealth, and incarceration status of the participants.
% Brief description of columns:
% race: Race of the respondent (hispanic, black or white),
% sex: Sex of the respondent (male or female),
% ever_jailed: A boolean variable that equals 1 if the respondent was jailed between 1985 to 1994,
% composite_wealth_1985: A composite variable creating by summing five wealth variables from the NLS data for the year 1985,
% composite_wealth_1990: A composite variable creating by summing five wealth variables from the NLS data for the year 1990,
% composite_wealth_1996: A composite variable creating by summing five wealth variables from the NLS data for the year 1996,

% Goal: What factors at the lowest end of the wealth distribution (10th percentile) do not significantly impact wealth when compared to higher ends of the wealth distribution (60th & 90th percentile)?
% """
% \end{minted}

% \subsection{Biology}

% Example 1: evolution\_freshwater\_fish $\rightarrow$ 3 $\rightarrow$ 0

% \begin{minted}[frame=lines,
% baselinestretch=1,
% bgcolor=Box3Color,
% fontsize=\footnotesize, breaklines, breakanywhere, breaksymbolleft={}, breaksymbolright={}]{python}
% task =
% f"""\

% Load all datasets using python using provided paths.
% Paths: ['../DiscoveryBench/evolution_freshwater_fish/body-size-evolution-in-south-american-freshwater-fishes.csv']

% Dataset name: body-size-evolution-in-south-american-freshwater-fishes.csv
% Dataset description: This dataset contains the drivers of speciation rates in South American freshwater fishes, employing an integrative approach that considers multiple biotic and abiotic factors.
% Brief description of columns:
% HYBAS_ID: Identifier for hydrological basins as defined by the HydroBASINS framework,
% long: The longitude for specific geographic locations used to map the presence/absence of species.,
% lat: The latitude for specific geographic locations used to map the presence/absence of species.,
% BAMM_speciation: This variable represents the average speciation rates of species within each sub-basin, as estimated by BAMM,
% BAMM_extinction: This variable represents the mean extinction rates for each tip of the phylogenetic tree, as estimated by the BAMM analysis,
% BAMM_NetDiv: This variable stands for the net diversification rates, which are calculated by subtracting the mean extinction rates from the mean
% speciation rates for each tip in the phylogenetic tree.,
% DR: Diversification Rate, a transformed measure of evolutionary distinctness to understand phylogenetic diversity. The lower the "ed" value, the higher the "DR" metric, indicating a species with less unique evolutionary history compared to others in the phylogeny.,
% BEL_evol: Rates of Body Elongation evolution,
% MBL_evol: Rates of Maximum Body Length evolution,
% OGP_evol: Rates of Oral Gape Position Evolution,
% RES_evol: Rates of Relative Eye Size Evolution,
% RML_evol: Rate of Relative Maxillary Length evolution,
% bio1: This variable represents the annual mean temperature.,
% bio12: This variable represents the annual mean precipitation.,
% runoff: This variable is used to represent the quantity of water from precipitation that flows over the land's surface and does not get absorbed into the ground. It is being extracted from a geographic database (HydroAtlas).,
% aet: Actual evapotranspiration, which is the sum of evaporation and plant transpiration from the Earth's land and ocean surface to the atmosphere.,
% Elevation: Mean elevation data indicating the height above sea level.,
% sgr: Stream gradient is a measure of the steepness or incline of a stream or river. It can affect water flow and sediment transport, which in turn can influence the habitat conditions for freshwater species. Higher stream gradients usually correspond to faster-moving
% water and can create different ecological conditions compared to lower gradients.,
% soil_div: It measures the diversity of soil types or conditions within each sub-basin studied. Soil diversity is computed using a dataset that includes eight variables related to substrate types and soil conditions.,
% area: The geographic area of a sub-basin, possibly in square meters, used as one of the predictors in the analysis.,
% diversity: Species diversity, which in ecological studies typically refers to the number of species and their relative abundances in a given area.,

% Goal: What entities show a weak but significant, positive relationship with a coefficient of 0.00003018?
% """
% \end{minted}

% Example 2: introduction\_pathways\_non-native\_plants $\rightarrow$ 4 $\rightarrow$ 2 

% \begin{minted}[frame=lines,
% baselinestretch=1,
% bgcolor=Box3Color,
% fontsize=\footnotesize, breaklines, breakanywhere, breaksymbolleft={}, breaksymbolright={}]{python}
% task =
% f"""\

% Load all datasets using python using provided paths.
% Paths: [../DiscoveryBench/introduction_pathways_non-native_plants/invaded_niche_pathways.csv, '../DiscoveryBench/introduction_pathways_non-native_plants/temporal_trends_contingency_table.csv']

% Dataset name: temporal-trends-contingency-table.csv
% Dataset description: Dataset contains temporal trends in the introduction pathways of non-native flora (plants) in the region of Catalonia.
% Brief description of columns:
% introduction.period: This column represents different time periods, related to when non-native plant species were introduced into the region.,
% pathway: This column represents different modes or routes through which non-native plant species were introduced, such as "AgriForest" (agriculture and forestry - plants introduced for cultivation to provide food or timber), "Gardening" (plants introduced for cultivation as ornamentals or for medicinal use), and "Unintentional" (plants introduced accidentally with the sowing of contaminated seed lots, global trade and tourism).,
% n: This column represents the frequency or count of non-native plant species introductions for each combination of introduction period and pathway.,

% Dataset name: invaded-niche-pathways.csv
% Dataset description: The data for the analysis of pathway-specific differences in the invaded niche.
% Brief description of columns:
% n.gard: The count or frequency of non-native plant species introduced through the "Gardening" pathway.,
% n.unint: The count or frequency of non-native plant species introduced through the "Unintentional" pathway,
% n.agfo: The count or frequency of non-native plant species introduced through the "AgriForest" (Agriculture and Forestry) pathway,
% n.total: The total count or frequency of non-native plant species across all introduction pathways.,
% habitat: A categorical variable representing the habitat type, selected from the ten most widespread habitat types in the Barcelona province,
% elevation: The elevation (in meters) of the sampled plot or location,
% cropland.1956.50m: This variable represent the percentage of cropland in a buffer of 50 meters (within a 50-meter radius) around the sampled plot or location for the year 1956,
% urban.1956.50m: This variable represent the percentage of urban land cover in a buffer (within a 50-meter radius) of 50 meters around the sampled plot or location for the year 1956,
% cropland.1993.50m: This variable represent the percentage of cropland in a buffer of 50 meters (within a 50-meter radius)  around the sampled plot or location for the year 1993,
% urban.1993.50m: This variable represent the percentage of urban land cover in a buffer of 50 meters (within a 50-meter radius) around the sampled plot or location for the year 1993,
% cropland.2009.50m: This variable represent the percentage of cropland in a buffer of 50 meters (within a 50-meter radius) around the sampled plot or location for the year 2009,
% ...
% ...
% urban.2009.1000m: This variable represents the percentage of urban land cover in a buffer of 1000 meters (within a 50-meter radius) around the sampled plot or location for the year 2009.,
% progressive.1956.2009.50m: This variable represents the historical landscape changes in a 50-meter buffer around the sampled plot or location between 1956 and 2009. It categorizes changes as progressive, indicating an increase in urban or cropland areas over this period.,
% regressive.1956.2009.50m: This variable represents the historical landscape changes in a 50-meter buffer around the sampled plot or location between 1956 and 2009. It categorizes changes as regressive, indicating a decrease in urban or cropland areas over this period.,
% no.changes.1956.2009.50m: his variable represents the historical landscape changes in a 50-meter buffer around the sampled plot or location between 1956 and 2009. It categorizes areas where there have been no significant changes in urban or cropland cover over this period.,
% progressive.1993.2009.50m: This variable represents the historical landscape changes in a 50-meter buffer around the sampled plot or location between 1993 and 2009. It categorizes changes as progressive, indicating an increase in urban or cropland areas over this period.,
% regressive.1993.2009.50m: This variable represents the historical landscape changes in a 50-meter buffer around the sampled plot or location between 1993 and 2009. It categorizes changes as regressive, indicating a decrease in urban or cropland areas over this period.,
% no.changes.1993.2009.50m: his variable represents the historical landscape changes in a 50-meter buffer around the sampled plot or location between 1993 and 2009. It categorizes areas where there have been no significant changes in urban or cropland cover over this period.,
% distance.stream: The distance (in meters) from the sampled plot or location to the nearest main stream or waterway.,
% distance.road: The distance (in meters) from the sampled plot or location to the nearest main road.,
% longitude: The geographic coordinates of the sampled plot or location.,
% latitude: The geographic coordinates of the sampled plot or location.,
% annual.temperature: The mean annual temperature of the sampled plot or location.,
% annual.rainfall: The annual precipitation or rainfall of the sampled plot or location.,
% annual.radiation: The mean annual solar radiation (in kJ/m²/day) of the sampled plot or location.,

% Goal: What is the relationship between introduction pathways and the success of non-native plants over time in Catalonia?
% """
% \end{minted}

% % Example 3: introduction\_pathways\_non-native\_plants||0||2

% % \begin{minted}[frame=lines,
% % baselinestretch=1,
% % bgcolor=Box3Color,
% % fontsize=\footnotesize, breaklines, breakanywhere, breaksymbolleft={}, breaksymbolright={}]{python}
% % task =
% % f"""\

% % Load all datasets using python using provided paths.
% % Paths: ['../DiscoveryBench/introduction_pathways_non-native_plants/temporal_trends_contingency_table.csv', './DiscoveryBench/introduction_pathways_non-native_plants/invaded_niche_pathways.csv']

% % Dataset name: temporal_trends_contingency_table.csv
% % Dataset description: Dataset contains temporal trends in the introduction pathways of non-native flora (plants) in the region of Catalonia.
% % Brief description of columns:
% % introduction.period: This column represents different time periods, related to when non-native plant species were introduced into the region.,
% % pathway: This column represents different modes or routes through which non-native plant species were introduced, such as "AgriForest" (agriculture and forestry - plants introduced for cultivation to provide food or timber), "Gardening" (plants introduced for cultivation as ornamentals or for medicinal use), and "Unintentional" (plants introduced accidentally with the sowing of contaminated seed lots, global trade and tourism).,
% % n: This column represents the frequency or count of non-native plant species introductions for each combination of introduction period and pathway.,
% % Dataset name: invaded_niche_pathways.csv
% % Dataset description: The dataset contains information about the different pathways in the invaded niche.
% % Brief description of columns:
% % n.gard: The count or frequency of non-native plant species introduced through the "Gardening" pathway.,
% % n.unint: The count or frequency of non-native plant species introduced through the "Unintentional" pathway,
% % n.agfo: The count or frequency of non-native plant species introduced through the "AgriForest" (Agriculture and Forestry) pathway,
% % n.total: The total count or frequency of non-native plant species across all introduction pathways.,
% % habitat: A categorical variable representing the habitat type, selected from the ten most widespread habitat types in the Barcelona province,
% % elevation: The elevation (in meters) of the sampled plot or location,
% % cropland.1956.50m: This variable represent the percentage of cropland in a buffer of 50 meters (within a 50-meter radius) around the sampled plot or location for the year 1956,
% % urban.1956.50m: This variable represent the percentage of urban land cover in a buffer (within a 50-meter radius) of 50 meters around the sampled plot or location for the year 1956,
% % cropland.1993.50m: This variable represent the percentage of cropland in a buffer of 50 meters (within a 50-meter radius)  around the sampled plot or location for the year 1993,
% % urban.1993.50m: This variable represent the percentage of urban land cover in a buffer of 50 meters (within a 50-meter radius) around the sampled plot or location for the year 1993,
% % cropland.2009.50m: This variable represent the percentage of cropland in a buffer of 50 meters (within a 50-meter radius) around the sampled plot or location for the year 2009,
% % urban.2009.50m: This variable represent the percentage of urban land cover in a buffer of 50 meters (within a 50-meter radius) around the sampled plot or location for the year 2009,
% % cropland.1956.500m: This variable represent the percentage of cropland in a buffer of 500 meters (within a 50-meter radius) around the sampled plot or location for the year 1956,
% % urban.1956.500m: This variable represent the percentage of urban land cover in a buffer of 50 meters (within a 50-meter radius) around the sampled plot or location for the year 1956,
% % cropland.1993.500m: This variable represent the percentage of cropland in a buffer of 500 meters (within a 50-meter radius) around the sampled plot or location for the year 1993,
% % urban.1993.500m: This variable represents the percentage of urban land cover in a buffer of 500 meters (within a 50-meter radius) around the sampled plot or location for the year 1993.,
% % cropland.2009.500m: This variable represents the percentage of cropland in a buffer of 500 meters (within a 50-meter radius) around the sampled plot or location for the year 2009.,
% % urban.2009.500m: This variable represents the percentage of urban land cover in a buffer of 500 meters (within a 50-meter radius) around the sampled plot or location for the year 2009.,
% % cropland.1956.1000m: This variable represents the percentage of cropland in a buffer of 1000 meters (within a 50-meter radius) around the sampled plot or location for the year 1956.,
% % urban.1956.1000m: This variable represents the percentage of urban land cover in a buffer of 1000 meters (within a 50-meter radius) around the sampled plot or location for the year 1956.,
% % cropland.1993.1000m: This variable represents the percentage of cropland in a buffer of 1000 meters (within a 50-meter radius) around the sampled plot or location for the year 1993.,
% % urban.1993.1000m: This variable represents the percentage of urban land cover in a buffer of 1000 meters (within a 50-meter radius) around the sampled plot or location for the year 1993.,
% % cropland.2009.1000m: This variable represents the percentage of cropland in a buffer of 1000 meters (within a 50-meter radius) around the sampled plot or location for the year 2009.,
% % urban.2009.1000m: This variable represents the percentage of urban land cover in a buffer of 1000 meters (within a 50-meter radius) around the sampled plot or location for the year 2009.,
% % progressive.1956.2009.50m: This variable represents the historical landscape changes in a 50-meter buffer around the sampled plot or location between 1956 and 2009. It categorizes changes as progressive, indicating an increase in urban or cropland areas over this period.,
% % regressive.1956.2009.50m: This variable represents the historical landscape changes in a 50-meter buffer around the sampled plot or location between 1956 and 2009. It categorizes changes as regressive, indicating a decrease in urban or cropland areas over this period.,
% % no.changes.1956.2009.50m: This variable represents the historical landscape changes in a 50-meter buffer around the sampled plot or location between 1956 and 2009. It categorizes areas where there have been no significant changes in urban or cropland cover over this period.,
% % progressive.1993.2009.50m: This variable represents the historical landscape changes in a 50-meter buffer around the sampled plot or location between 1993 and 2009. It categorizes changes as progressive, indicating an increase in urban or cropland areas over this period.,
% % regressive.1993.2009.50m: This variable represents the historical landscape changes in a 50-meter buffer around the sampled plot or location between 1993 and 2009. It categorizes changes as regressive, indicating a decrease in urban or cropland areas over this period.,
% % no.changes.1993.2009.50m: This variable represents the historical landscape changes in a 50-meter buffer around the sampled plot or location between 1993 and 2009. It categorizes areas where there have been no significant changes in urban or cropland cover over this period.,
% % distance.stream: The distance (in meters) from the sampled plot or location to the nearest main stream or waterway.,
% % distance.road: The distance (in meters) from the sampled plot or location to the nearest main road.,
% % longitude: The geographic coordinates of the sampled plot or location.,
% % latitude: The geographic coordinates of the sampled plot or location.,
% % annual.temperature: The mean annual temperature of the sampled plot or location.,
% % annual.rainfall: The annual precipitation or rainfall of the sampled plot or location.,
% % annual.radiation: The mean annual solar radiation (in kJ/m²/day) of the sampled plot or location.,

% % Goal: In what way has gardening impacted agriculture as a contributor to non-native flora over the past millennium?

% % 
% % """
% % \end{minted}


% \subsection{Economics}

% % Example 1: worldbank\_education\_gdp||2||0

% % \begin{minted}[frame=lines,
% % baselinestretch=1,
% % bgcolor=Box3Color,
% % fontsize=\footnotesize, breaklines, breaksymbolleft={}, breaksymbolright={}]{python}
% % task =
% % f"""\

% % Load all datasets using python using provided paths.
% % Paths: ['../DiscoveryBench/worldbank_education_gdp/worldbank_education_gdp.csv']
% % Dataset name: worldbank_education_gdp.csv
% % Dataset description: This is the panel data created based on the information provided in the paper. The dataset contains information of two groups of countries- Sub-Saharan Africa and Lower Middle Income Countries (LMC) from 1975 to 2015, sourced from the World Development Indicators (WDI) database. It focuses on the relationship between government expenditure on education (% of total expenditure) and per capita GDP (in 2010 US$), using variables like total labor force, gross primary and secondary enrollment, and exports (annual % growth) as key indicators.
% % Brief description of columns:
% % Country Group: The name given to the group of countries- Sub-Saharan Africa and Lower middle income,
% % Country Code: The code name assigned to each group of countries: Sub-Saharan Africa-SSA and Lower Middle Income Countries-LMC,
% % Series Name: The name of the indicator or variable being measured. Description of the different indicators in the series is as follows:-GNI per capita (constant 2015 US dollar) - Gross National Income (GNI) per capita adjusted for inflation to 2015 US dollars. Adjusted savings: education expenditure (percentage of GNI) - The percentage of GNI that is saved for future education expenditures. Exports of goods and services (annual percentage growth) - The annual percentage growth rate of exports of goods and services. School enrollment, primary (percentage gross) - The gross enrollment ratio for primary school, representing the percentage of children of official primary school age who are enrolled in primary school. School enrollment, secondary (percentage gross) - The gross enrollment ratio for secondary school, representing the percentage of children of official secondary school age who are enrolled in secondary school. Labor force participation rate, total (percentage of total population ages 15+) (modeled ILO estimate) - The percentage of the total population ages 15 and older that is economically active.,
% % Series Code: The code representing the indicator or variable,
% % 1975 [YR1975]: The value of each indicator or variable for the year 1975,
% % 1976 [YR1976]: The value of each indicator or variable for the year 1976,
% % ...
% % ...
% % ...
% % 2014 [YR2014]: The value of each indicator or variable for the year 2014,
% % 2015 [YR2015]: The value of each indicator or variable for the year 2015,

% % Goal: How does the effect of increasing education expenditure on per capita GDP compare between developing countries and countries in Sub-Saharan Africa?

% % """
% % \end{minted}

% Example 1: worldbank\_education\_gdp\_indicators $\rightarrow$ 4 $\rightarrow$ 0

% \begin{minted}[frame=lines,
% baselinestretch=1,
% bgcolor=Box3Color,
% fontsize=\footnotesize, breaklines, breakanywhere, breaksymbolleft={}, breaksymbolright={}]{python}
% task =
% f"""\

% Load all datasets using python using provided paths.
% Paths: Paths: ['../DiscoveryBench/worldbank_education_gdp_indicators/Adjusted_savings_education_expenditure_percentage_of_GNI.csv', '../DiscoveryBench/worldbank_education_gdp_indicators/Exports_of_goods_and_services_annual_percentage_growth.csv', '../DiscoveryBench/worldbank_education_gdp_indicators/GNI_per_capita_constant_2015_USdollar.csv', '../DiscoveryBench/worldbank_education_gdp_indicators/Labor_force_participation_rate_total_percentage_of_total_population_ages_15+_modeled_ILO_estimate.csv', '../DiscoveryBench/worldbank_education_gdp_indicators/School_enrollment_primary_percentage_gross.csv', '../DiscoveryBench/worldbank_education_gdp_indicators/School_enrollment_secondary_percentage_gross.csv']


% Dataset name: Adjusted_savings_education_expenditure_percentage_of_GNI.csv
% Dataset description: This dataset contains information on adjusted savings: education expenditure (percentage of gni) for Sub-Saharan Africa and Lower Middle Income Countries (LMC) from 1975 to 2015.
% Brief description of columns:
% Country Group: The name given to the group of countries- Sub-Saharan Africa and Lower middle income,
% Country Code: The code name assigned to each group of countries: Sub-Saharan Africa-SSA and Lower Middle Income Countries-LMC,
% 1975 [YR1975]: The value of indicator for the year 1975,
% 1976 [YR1976]: The value of indicator for the year 1976,
% ...
% ...
% 2013 [YR2013]: The value of indicator for the year 2013,
% 2014 [YR2014]: The value of indicator for the year 2014,
% 2015 [YR2015]: The value of indicator for the year 2015,

% Dataset name: Exports_of_goods_and_services_annual_percentage_growth.csv
% Dataset description: This dataset contains information on exports of goods and services (annual percentage growth) for Sub-Saharan Africa and Lower Middle Income Countries (LMC) from 1975 to 2015.
% Brief description of columns:
% Country Group: The name given to the group of countries- Sub-Saharan Africa and Lower middle income,
% Country Code: The code name assigned to each group of countries: Sub-Saharan Africa-SSA and Lower Middle Income Countries-LMC,
% 1975 [YR1975]: The value of indicator for the year 1975,
% 1976 [YR1976]: The value of indicator for the year 1976,
% ...
% ...
% 2014 [YR2014]: The value of indicator for the year 2014,
% 2015 [YR2015]: The value of indicator for the year 2015,

% Dataset name: GNI_per_capita_constant_2015_USdollar.csv
% Dataset description: This dataset contains information on gni per capita (constant 2015 usdollar) for Sub-Saharan Africa and Lower Middle Income Countries (LMC) from 1975 to 2015.
% Brief description of columns:
% Country Group: The name given to the group of countries- Sub-Saharan Africa and Lower middle income,
% Country Code: The code name assigned to each group of countries: Sub-Saharan Africa-SSA and Lower Middle Income Countries-LMC,

% 1975 [YR1975]: The value of indicator for the year 1975,
% 1976 [YR1976]: The value of indicator for the year 1976,
% ...
% ...
% 2013 [YR2013]: The value of indicator for the year 2013,
% 2014 [YR2014]: The value of indicator for the year 2014,
% 2015 [YR2015]: The value of indicator for the year 2015,

% Dataset name: Labor_force_participation_rate_total_percentage_of_total_population_ages_15+_modeled_ILO_estimate.csv
% Dataset description: This dataset contains information on labor force participation rate, total (percentage of total population ages 15+) (modeled ilo estimate) for Sub-Saharan Africa and Lower Middle Income Countries (LMC) from 1975 to 2015.
% Brief description of columns:
% Country Group: The name given to the group of countries- Sub-Saharan Africa and Lower middle income,
% Country Code: The code name assigned to each group of countries: Sub-Saharan Africa-SSA and Lower Middle Income Countries-LMC,
% 1975 [YR1975]: The value of indicator for the year 1975,
% 1976 [YR1976]: The value of indicator for the year 1976,
% ...
% ...
% 2014 [YR2014]: The value of indicator for the year 2014,
% 2015 [YR2015]: The value of indicator for the year 2015,

% Dataset name: School_enrollment_primary_percentage_gross.csv
% Dataset description: This dataset contains information on school enrollment, primary (percentage gross) for Sub-Saharan Africa and Lower Middle Income Countries (LMC) from 1975 to 2015.
% Brief description of columns:
% Country Group: The name given to the group of countries- Sub-Saharan Africa and Lower middle income,
% Country Code: The code name assigned to each group of countries: Sub-Saharan Africa-SSA and Lower Middle Income Countries-LMC,
% 1975 [YR1975]: The value of indicator for the year 1975,
% 1976 [YR1976]: The value of indicator for the year 1976,
% ...
% ...
% 2015 [YR2015]: The value of indicator for the year 2015,

% Dataset name: School_enrollment_secondary_percentage_gross.csv
% Dataset description: This dataset contains information on school enrollment, secondary (percentage gross) for Sub-Saharan Africa and Lower Middle Income Countries (LMC) from 1975 to 2015.
% Brief description of columns:
% Country Group: The name given to the group of countries- Sub-Saharan Africa and Lower middle income,
% Country Code: The code name assigned to each group of countries: Sub-Saharan Africa-SSA and Lower Middle Income Countries-LMC,
% 1975 [YR1975]: The value of indicator for the year 1975,
% 1976 [YR1976]: The value of indicator for the year 1976,
% ...
% ...
% 2014 [YR2014]: The value of indicator for the year 2014,
% 2015 [YR2015]: The value of indicator for the year 2015,

% Goal: How do labor productivity and education levels relate to economic output, particularly in terms of export growth?

% """
% \end{minted}

% Example 2:  immigration\_offshoring\_effect\_on\_employment $\rightarrow$ 1 $\rightarrow$ 1

% \begin{minted}[frame=lines,
% baselinestretch=1,
% bgcolor=Box3Color,
% fontsize=\footnotesize, breaklines, breakanywhere, breaksymbolleft={}, breaksymbolright={}]{python}
% task =
% f"""\

% Load all datasets using python using provided paths.
% Paths: ['../DiscoveryBench/immigration_offshoring_effect_on_employment/offshoring_iv_mar2.dta','../DiscoveryBench/immigration_offshoring_effect_on_employment/immi_popimputed_00_07.dta']

% Dataset name: offshoring_iv_mar2.dta
% Dataset description: This dataset contains measures aimed at capturing exogenous variation in the ease or costs of offshoring across industries and years. It is constructed using variation in offshoring across countries to the U.S., interacting with each industry's initial distribution of offshoring across those countries. This offshoring measure is intended to be used as an explanatory variable when examining impacts on domestic employment patterns.
% Brief description of columns:
% year: The year of the observation,
% beaind: Beaurau of Economic (BEA) Industry Code,
% iv_offshoring_1: The key instrumental variable capturing exogenous variation in the ease/costs of offshoring for that industry-year.,
% Dataset name: immi_popimputed_00_07.dta
% Dataset description: This dataset provides imputed measures of immigrant employment aimed at capturing exogenous variation in immigration costs and push-factors across industries and years
% Brief description of columns:
% year: The year of the observation,
% beaind: Beaurau of Economic (BEA) Industry Code,
% share_immi_imputed: This column represents the imputed share or proportion of immigrant employment within total employment for each industry-year observation. The share of immigrant employment indicates the proportion of total employment within each industry that is comprised of immigrant workers.,
% empl_immi_imputed: This column represents the imputed level of immigrant employment for each industry-year observation. This refers to the estimated number of immigrant workers employed within each industry-year.,

% Goal: How does increased ease of immigration impact the share of native employment?

% """
% \end{minted}

% \subsection{Engineering}

% requirements\_engineering\_for\_ML\_enabled\_systems $\rightarrow$ 14 $\rightarrow$ 0

% \begin{minted}[frame=lines,
% baselinestretch=1,
% bgcolor=Box3Color,
% fontsize=\footnotesize, breaklines, breakanywhere, breaksymbolleft={}, breaksymbolright={}]{python}
% task =
% f"""\

% Load all datasets using python using provided paths.
% Paths: ['../DiscoveryBench/requirements_engineering_for_ML_enabled_systems/requirements_engineering_for_ML-enabled_systems.csv']

% Dataset name: requirements_engineering_for_ML-enabled_systems.csv
% Dataset description: Survey responses detailing the roles, techniques, and documentation practices associated with requirements in ML-enabled system projects.
% Brief description of columns:
% ID: The unique identifier for each respondent.,
% Status: The current status of the respondent,
% Duration: The duration of the respondent's involvement,
% D1_Undergraduation: Undergraduate (e.g., Computer Science, Statistics),
% D1_Specialization: Specialization (e.g., Data Science specialization, Project Management specialization),
% D1_Master: Master (e.g., M.Sc. in Computer Science, M.Sc. in Economics),
% D1_Phd: Ph.D. (e.g., Ph.D. in Computer Science, Ph.D. in Mathematics),
% D1_Courses: Professional ML Certifications/Courses (e.g., Google Professional ML Engineer Certification, Coursera/Udacity course on ML),
% D1_Others: Other course specified by respondent,
% D2_Country: Country in which the respondent is currently working,
% D3_Company_Size: Size of the organization the respondent currently work for (1-10 employees, 11-50 employees ... more than 2000 employees),
% D4_Role: Role that best describes the respondent's current activities within the company (Project Lead/ Project Manager, business Analyst, Requirements Engineer, Solution Architect, Data Scientist, Developer, Test Manager / Tester),
% D4_Role_Others: Other role specified by respondent,
% D5_Software_Experience: Years of experience in working with the development of software based products,
% D6_ML_Experience: Years of Experience in developing ML-enabled systems,
% D7_Total_ML_Projects: Number of ML-enabled system projects that the respondent participated in,
% D8_ML_Production: Number of ML-enabled system projects that the respondent participated in that actually got deployed,
% D9_ML_Project_Team_Size: The Team size of the ML-enabled system projects that the respondent participated in,
% D10_ML_Management_Framework_None: Participant responded with None as the response for project management framework applied in the participated ML-enabled systems project,      
% D10_ML_Management_Framework_CRISP-DM: Participant responded with CRISP-DM as the response for project management framework applied in the participated ML-enabled systems project,
% D10_ML_Management_Framework_Kanban: Participant responded with Kanban as the response for project management framework applied in the participated ML-enabled systems project,  
% D10_ML_Management_Framework_Lean: Participant responded with Lean as the response for project management framework applied in the participated ML-enabled systems project,      
% D10_ML_Management_Framework_RUP: Participant responded with RUP as the response for project management framework applied in the participated ML-enabled systems project,        
% D10_ML_Management_Framework_SAFe: Participant responded with SAFe as the response for project management framework applied in the participated ML-enabled systems project,      
% D10_ML_Management_Framework_Scrum: Participant responded with Scrum as the response for project management framework applied in the participated ML-enabled systems project,    
% D10_ML_Management_Framework_Others: Participant responded with a different framework as the response for project management framework applied in the participated ML-enabled systems project,
% D10_ML_Management_Framework_Others_Free: Name of the other framework for project management framework applied in the participated ML-enabled systems project,

% D11_Agile_Development: The agility of the development of the respondent in the ML-enabled systems projects that the respondent participated in,
% D12_ML_Project_Context_Banking: Banking was the domain of application of the ML-enabled systems project that the respondent participated in,
% D12_ML_Project_Context_Defense: Defense was the domain of application of the ML-enabled systems project that the respondent participated in,
% D12_ML_Project_Context_Education: Education was the domain of application of the ML-enabled systems project that the respondent participated in,
% D12_ML_Project_Context_Embedded: Embedded systems in Automotive or Avionics was the domain of application of the ML-enabled systems project that the respondent participated in,
% D12_ML_Project_Context_Entertainment: Entertainment was the domain of application of the ML-enabled systems project that the respondent participated in,
% D12_ML_Project_Context_Healthcare: Healthcare was the domain of application of the ML-enabled systems project that the respondent participated in,
% D12_ML_Project_Context_Insurance: Insurance was the domain of application of the ML-enabled systems project that the respondent participated in,
% D12_ML_Project_Context_Logistics: Logistics was the domain of application of the ML-enabled systems project that the respondent participated in,
% D12_ML_Project_Context_Oil: Oil & Gas was the domain of application of the ML-enabled systems project that the respondent participated in,
% D12_ML_Project_Context_Sales: Sales/E-commerce was the domain of application of the ML-enabled systems project that the respondent participated in,
% D12_ML_Project_Context_Telecom: Telecommunication was the domain of application of the ML-enabled systems project that the respondent participated in,
% D12_ML_Project_Context_Others: Respondent specified some other domain of application of the ML-enabled systems project that the respondent participated in,
% D12_ML_Project_Context_Others_Free: Respondent's domain of application of the ML-enabled systems project that the respondent participated in,
% D13_ML_Programming_Language_C: C language was in the list of general languages that composed the ML-enabled system projects (including eventually Non-ML related parts),        
% D13_ML_Programming_Language_Java: Java language was in the list of general languages that composed the ML-enabled system projects (including eventually Non-ML related parts),  
% D13_ML_Programming_Language_Javascript: Javascript language was in the list of general languages that composed the ML-enabled system projects (including eventually Non-ML related parts),
% D13_ML_Programming_Language_Julia: Julia language was in the list of general languages that composed the ML-enabled system projects (including eventually Non-ML related parts),
% D13_ML_Programming_Language_MatLab: MatLab language was in the list of general languages that composed the ML-enabled system projects (including eventually Non-ML related parts),
% ...
% ...
% D14_ML_Purpose_Association: Association was the main purpose of the ML-enabled system projects the respondent participated in,
% D14_ML_Purpose_Association_Free: The typical purposes that were addressed using association in the project,
% D14_ML_Purpose_Clustering: Clustering was the main purpose of the ML-enabled system projects the respondent participated in,
% D14_ML_Purpose_Clustering_Free: The typical purposes that were addressed using clustering in the project,
% D14_ML_Purpose_Others: ML-enabled system project had some other purpose,
% D14_ML_Purpose_Others_Free: The other purposes that were addressed in the project,
% D15_ML_Algorithms_Apriori: Apriori algorithm was employed in the ML-enabled system project that the respondent participated in,
% D15_ML_Algorithms_Bayesian: Bayesian algorithm was employed in the ML-enabled system project that the respondent participated in,
% D15_ML_Algorithms_DBSCAN: DBSCAN algorithm was employed in the ML-enabled system project that the respondent participated in,
% D15_ML_Algorithms_Decision_Tree: Decision Tree algorithm was employed in the ML-enabled system project that the respondent participated in,
% D15_ML_Algorithms_Ensembles: Ensemble  (e.g. Random Forests, XGBoost) Algorithm was employed in the ML-enabled system project that the respondent participated in,
% ...
% ...
% D15_ML_Algorithms_Naive_Bayes: Naive Bayes was employed in the ML-enabled system project that the respondent participated in,
% D15_ML_Algorithms_Neural_Networks: Neural Networks were employed in the ML-enabled system project that the respondent participated in,
% D15_ML_Algorithms_SVM: Support Vector Machines was employed in the ML-enabled system project that the respondent participated in,
% D15_ML_Algorithms_Others: Some other algorithm was employed in the ML-enabled system project that the respondent participated in,
% D15_ML_Algorithms_Others_Free: The name of the different algorithm that was employed in the ML-enabled system project that the respondent participated in,
% Q1_ML_Life_Cycle_Importance_Problem_Understanding: The level of relevance of Problem Understanding and Requirements with regard to project success. One of the following: Not Relevant at All, Low Relevance, Neutral, High Relevance, Extremely Relevant, I don't know,
% Q1_ML_Life_Cycle_Importance_Data_Collection: The level of relevance of Data Collection with regard to project success. One of the following: Not Relevant at All, Low Relevance, Neutral, High Relevance, Extremely Relevant, I don't know,
% Q1_ML_Life_Cycle_Importance_Data_Pre-Processing: The level of relevance of Data Pre-Processing with regard to project success. One of the following: Not Relevant at All, Low Relevance, Neutral, High Relevance, Extremely Relevant, I don't know,
% Q1_ML_Life_Cycle_Importance_Model_Creation: The level of relevance of Model Creation with regard to project success. One of the following: Not Relevant at All, Low Relevance, Neutral, High Relevance, Extremely Relevant, I don't know,
% Q1_ML_Life_Cycle_Importance_Model_Evaluation: The level of relevance of Model Evaluation with regard to project success. One of the following: Not Relevant at All, Low Relevance, Neutral, High Relevance, Extremely Relevant, I don't know,
% Q1_ML_Life_Cycle_Importance_Model_Deployment: The level of relevance of Model Deployment with regard to project success. One of the following: Not Relevant at All, Low Relevance, Neutral, High Relevance, Extremely Relevant, I don't know,
% Q1_ML_Life_Cycle_Importance_Model_Monitoring: The level of relevance of Model Monitoring with regard to project success. One of the following: Not Relevant at All, Low Relevance, Neutral, High Relevance, Extremely Relevant, I don't know,
% Q2_ML_Life_Cycle_Difficulty_Problem_Understanding: Difficulty level of Problem Understanding and Requirements stage in ML Life Cycle. One of the following: Very Easy, Easy, Neutral, Complex, Very Complex, I don't know,
% Q2_ML_Life_Cycle_Difficulty_Data_Collection: Difficulty level of Data Collection stage in ML Life Cycle. One of the following: Very Easy, Easy, Neutral, Complex, Very Complex, I don't know,
% Q2_ML_Life_Cycle_Difficulty_Data_Pre-Processing: Difficulty level of Data Pre-Processing stage in ML Life Cycle. One of the following: Very Easy, Easy, Neutral, Complex, Very Complex, I don't know,
% Q2_ML_Life_Cycle_Difficulty_Model_Creation: Difficulty level of Model Creation stage in ML Life Cycle. One of the following: Very Easy, Easy, Neutral, Complex, Very Complex, I don't know,
% Q2_ML_Life_Cycle_Difficulty_Model_Evaluation: Difficulty level of Model Evaluation stage in ML Life Cycle. One of the following: Very Easy, Easy, Neutral, Complex, Very Complex, I don't know,
% Q2_ML_Life_Cycle_Difficulty_Model_Deployment: Difficulty level of Model Deployment stage in ML Life Cycle. One of the following: Very Easy, Easy, Neutral, Complex, Very Complex, I don't know,
% Q2_ML_Life_Cycle_Difficulty_Model_Monitoring: Difficulty level of Model Monitoring stage in ML Life Cycle. One of the following: Very Easy, Easy, Neutral, Complex, Very Complex, I don't know,
% Q3_ML_Life_Cycle_Effort_Problem_Understanding: The proportion of effort spent in the ML life cycle stage for Problem Understanding,
% Q3_ML_Life_Cycle_Effort_Data_Collection: The proportion of effort spent in the ML life cycle stage for Data Collection,
% Q3_ML_Life_Cycle_Effort_Data_Pre-Processing: The proportion of effort spent in the ML life cycle stage for Data Pre-Processing,
% Q3_ML_Life_Cycle_Effort_Model_Creation: The proportion of effort spent in the ML life cycle stage for Model Creation,
% Q3_ML_Life_Cycle_Effort_Model_Evaluation: The proportion of effort spent in the ML life cycle stage for Model Evaluation,
% Q3_ML_Life_Cycle_Effort_Model_Deployment: The proportion of effort spent in the ML life cycle stage for Model Deployment,
% Q3_ML_Life_Cycle_Effort_Model_Monitoring: The proportion of effort spent in the ML life cycle stage for Model Monitoring,
% Q4_ML_Life_Cycle_Main_Problems_Problem_Understanding_Free_First: The first main problem faced in Problem Understanding phase in the ML life cycle stage,
% Q4_ML_Life_Cycle_Main_Problems_Problem_Understanding_Free_Second: The second main problem faced in the Problem Understanding phase of the ML life cycle,
% ...
% ...
% ...
% Q16_Model_Monitor_Aspects_Input_And_Output: Importance of monitoring inputs and outputs of models in the respondent's organization,
% Q16_Model_Monitor_Aspects_Interpretability_Output: Importance of monitoring the interpretability of model outputs in the respondent's organization,
% Q16_Model_Monitor_Aspects_Output_And_Decisions: Importance of monitoring outputs and decisions of models in the respondent's organization,
% Q16_Model_Monitor_Aspects_Fairness: Importance of monitoring fairness of models in the respondent's organization,
% Q16_Model_Monitor_Aspects_Others: Importance of monitoring other aspects of models specified by the respondent,
% Q16_Model_Monitor_Aspects_Others_Free: Free text response for other aspects of model monitoring specified by the respondent,
% Q17_Automated_Machine_Learning_Tools_Yes_No: Yes or No response indicating if the respondent uses automated machine learning tools,
% Q17_Automated_Machine_Learning_Tools_Yes_Free: Free text response if the respondent uses automated machine learning tools,
% Origin: Origin of the respondent,

% Goal: What are the percentages of respondents and the 95% Confidence Interval of the percentage after bootstrapping for statistical significance for each of the following tasks: 1) aligning requirements data, 2) changing requirements, 3) managing conflicts, and 4) selecting metrics where they are considered significantly difficult when defining requirements for ML-enabled systems?

% """
% \end{minted}


% \subsection{Meta-science}

% meta\_regression\_raw $\rightarrow$ 0 $\rightarrow$ 3


% \begin{minted}[frame=lines,
% baselinestretch=1,
% bgcolor=Box3Color,
% fontsize=\footnotesize, breaklines, breakanywhere, breaksymbolleft={}, breaksymbolright={}]{python}
% task =
% f"""\

% Load all datasets using python using provided paths.
% Paths: ['../DiscoveryBench/meta_regression_raw/meta-regression_replication_success_data_heterogeneity_in_replication_projects.csv', '../DiscoveryBench/meta_regression_raw/meta-regression_study_data_heterogeneity_in_replication_projects.csv']

% Dataset name: meta-regression_study_data_heterogeneity_in_replication_projects.csv
% Dataset description: Dataset contains information about original & replication studies. Original & replication specific columns may be appended by o & r
% Brief description of columns:
% id: Unique id for each O/R pair,
% title: Title of the research study,
% authors.o: Names of Original paper's authors,
% pub_year: Year of Publication of the study,
% journal: Journal in which the study was published,
% volume: Volume Number of the journal,
% issue: Issue Number of the journal,
% discipline: Discipline of original paper. One of the following: Social, Cognitive or Economics,
% length: Number of pages of original paper,
% citations: Number of citations of original paper,
% effect_size.o: Standardized effect size of original paper,
% p_value.o: P-value of original paper,
% n.o: Sample size of original paper,
% effect_type: Type of effect tested. One of the following: main effect, correlation, interaction,
% effect_size.r: Standardized effect size of replication,
% p_value.r: P-value of replication,
% n_planned.r: Planned sample size of replication,
% n.r: Sample size of replication,
% power.o: Post hoc power based on original effect size,
% power.r: Post hoc power based on replication effect size,
% power_planned.r: Planned power of the replication based on planned N and original ES,
% experiment_country.o: Country where original experiment was conducted,
% experiment_country.r: Country where replication is to be conducted,
% experiment_language.o: Language used with subjects in original experiment (English, German, Dutch, Polish, Hebrew, French, Italian, Arabic, Spanish, Korean),  
% experiment_language.r: Language to be used with subjects in replication (English, Polish, German, Dutch, Italian, Portuguese, Malay, Turkish, Czech, Arabic, Spanish),
% online.o: If the original experiment was conducted online (1: yes, 0: no),
% online.r: If the replication was conducted online (1: yes, 0: no),
% compensation.o: Compensation in original experiment (credit, cash, nothing, mixed),
% compensation.r: Compensation in replication (credit, cash, nothing, mixed),
% subjects.o: Type of subjects used in original experiment (students, online, anyone, community),
% subjects.r: Type of subjects used in replication (students, online, anyone, community),
% endprice: Final market price in prediction market,
% transactions: Number of transactions in prediction market,
% trading_volume: Total volume of traded stocks in prediction market,
% replicated: Binary outcome variable; study is replicated if p <= 0.05 and effect goes in the same direction as the original,
% project: The replication project that the study was in (ml1: Many Labs 1, ml3: Many Labs 3, rpp: Psychology, ee: Experimental Economics),
% relative_es: The continuous outcome variable; the standardized replication effect size to the original effect | relative effect size = (replication effect size / original effect size),
% n_authors.o: Number of authors in original study,
% n_authors.r: Number of authors in replication,
% author_citations_avg.o: Average number of citations of authors in original study,
% author_citations_max.o: Number of citations of the author in original study with the highest citation count,
% authors_male.o: Ratio of male authors in original study,
% seniority.o: Most senior author in the original paper (Professor, Associate Professor, Assistant, Researcher, Assistant Professor),
% author_citations_avg.r: Average number of citations of authors in replication study,
% author_citations_max.r: Number of citations of the author in original study with the highest citation count,
% authors_male.r: Ratio of male authors in replication,
% seniority.r: Most senior author in the original paper (Professor, Associate Professor, Assistant, Researcher, Assistant Professor),
% aggregated: Aggregated column,
% lab_id: Unique id for each replication lab,
% es_80power: Standardized effect size required in replication to achieve 80% power,
% same_country: Original study and replication are in the same country,
% same_language: Original study and replication are in the same language,
% same_online: Original study and replication are both conducted online,
% same_subjects: Original study and replication use same type of subjects,
% us_lab.o: Original experiment lab in the US,
% us_lab.r: Replication experiment lab in the US,

% Dataset name: meta-regression_replication_success_data_heterogeneity_in_replication_projects.csv
% Dataset description: Data from four large-scale replication projects
% Brief description of columns:
% study: Study identifier, usually names of authors from original study,
% project: Name of replication project,
% ro: Effect estimate of original study on correlation scale,
% rr: Effect estimate of replication study on correlation scale,
% fiso: Effect estimate of original study transformed to Fisher-z scale,
% fisr: Effect estimate of replication study transformed to Fisher-z scale,
% se_fiso: Standard error of Fisher-z transformed effect estimate of original study,
% se_fisr: Standard error of Fisher-z transformed effect estimate of replication study,
% po: Two-sided p-value from significance test of effect estimate from original study,
% pr: Two-sided p-value from significance test of effect estimate from replication study,
% po1: One-sided p-value from significance test of effect estimate from original study (in the direction of the original effect estimate),
% pr1: One-sided p-value from significance test of effect estimate from replication study (in the direction of the original effect estimate),
% pm_belief: Peer belief about whether replication effect estimate will achieve statistical significance elicited through prediction market (only available for EERP and SSRP),
% no: Sample size in original study,
% nr: Sample size in replication study,

% Goal: Which factor in Experimental Economics has a value of 0.57 on the Fisher-z scale in original studies compared to 0.31 in replication studies?

% """
% \end{minted}

% % Example 2: meta\_regression\_raw||19||0

% % \begin{minted}[frame=lines,
% % baselinestretch=1,
% % bgcolor=Box3Color,
% % fontsize=\footnotesize, breaklines, breakanywhere, breaksymbolleft={}, breaksymbolright={}]{python}
% % task =
% % f"""\

% % Load all datasets using python using provided paths.
% % Paths: ['../DiscoveryBench/meta_regression_raw/meta-regression_replication_success_data_heterogeneity_in_replication_projects.csv', '../DiscoveryBench/meta_regression_raw/meta-regression_study_data_heterogeneity_in_replication_projects.csv']

% % Dataset name: meta-regression_study_data_heterogeneity_in_replication_projects.csv
% % Dataset description: Dataset contains information about original & replication studies. Original & replication specific columns may be appended by o & r
% % Brief description of columns:
% % id: Unique id for each O/R pair,
% % title: Title of the research study,
% % authors.o: Names of Original paper's authors,
% % pub_year: Year of Publication of the study,
% % journal: Journal in which the study was published,
% % volume: Volume Number of the journal,
% % issue: Issue Number of the journal,
% % discipline: Discipline of original paper. One of the following: Social, Cognitive or Economics,
% % length: Number of pages of original paper,
% % citations: Number of citations of original paper,
% % effect_size.o: Standardized effect size of original paper,
% % p_value.o: P-value of original paper,
% % n.o: Sample size of original paper,
% % effect_type: Type of effect tested. One of the following: main effect, correlation, interaction,
% % effect_size.r: Standardized effect size of replication,
% % p_value.r: P-value of replication,
% % n_planned.r: Planned sample size of replication,
% % n.r: Sample size of replication,
% % power.o: Post hoc power based on original effect size,
% % power.r: Post hoc power based on replication effect size,
% % power_planned.r: Planned power of the replication based on planned N and original ES,
% % experiment_country.o: Country where original experiment was conducted,
% % experiment_country.r: Country where replication is to be conducted,
% % experiment_language.o: Language used with subjects in original experiment (English, German, Dutch, Polish, Hebrew, French, Italian, Arabic, Spanish, Korean),
% % experiment_language.r: Language to be used with subjects in replication (English, Polish, German, Dutch, Italian, Portuguese, Malay, Turkish, Czech, Arabic, Spanish),
% % online.o: If the original experiment was conducted online (1: yes, 0: no),
% % online.r: If the replication was conducted online (1: yes, 0: no),
% % compensation.o: Compensation in original experiment (credit, cash, nothing, mixed),
% % compensation.r: Compensation in replication (credit, cash, nothing, mixed),
% % subjects.o: Type of subjects used in original experiment (students, online, anyone, community),
% % subjects.r: Type of subjects used in replication (students, online, anyone, community),
% % endprice: Final market price in prediction market,
% % transactions: Number of transactions in prediction market,
% % trading_volume: Total volume of traded stocks in prediction market,
% % replicated: Binary outcome variable; study is replicated if p <= 0.05 and effect goes in the same direction as the original,
% % project: The replication project that the study was in (ml1: Many Labs 1, ml3: Many Labs 3, rpp: Psychology, ee: Experimental Economics), 
% % relative_es: The continuous outcome variable; the standardized replication effect size to the original effect | relative effect size = (replication effect size / original effect size),
% % n_authors.o: Number of authors in original study,
% % n_authors.r: Number of authors in replication,
% % author_citations_avg.o: Average number of citations of authors in original study,
% % author_citations_max.o: Number of citations of the author in original study with the highest citation count,
% % authors_male.o: Ratio of male authors in original study,
% % seniority.o: Most senior author in the original paper (Professor, Associate Professor, Assistant, Researcher, Assistant Professor),       
% % author_citations_avg.r: Average number of citations of authors in replication study,
% % author_citations_max.r: Number of citations of the author in original study with the highest citation count,
% % authors_male.r: Ratio of male authors in replication,
% % seniority.r: Most senior author in the original paper (Professor, Associate Professor, Assistant, Researcher, Assistant Professor),       
% % aggregated: Aggregated column,
% % lab_id: Unique id for each replication lab,
% % es_80power: Standardized effect size required in replication to achieve 80% power,
% % same_country: Original study and replication are in the same country,
% % same_language: Original study and replication are in the same language,
% % same_online: Original study and replication are both conducted online,
% % same_subjects: Original study and replication use same type of subjects,
% % us_lab.o: Original experiment lab in the US,
% % us_lab.r: Replication experiment lab in the US,
% % drop: Drop column,
% % Dataset name: meta-regression_replication_success_data_heterogeneity_in_replication_projects.csv
% % Dataset description: Data from four large-scale replication projects
% % Brief description of columns:
% % study: Study identifier, usually names of authors from original study,
% % project: Name of replication project,
% % ro: Effect estimate of original study on correlation scale,
% % rr: Effect estimate of replication study on correlation scale,
% % fiso: Effect estimate of original study transformed to Fisher-z scale,
% % fisr: Effect estimate of replication study transformed to Fisher-z scale,
% % se_fiso: Standard error of Fisher-z transformed effect estimate of original study,
% % se_fisr: Standard error of Fisher-z transformed effect estimate of replication study,
% % po: Two-sided p-value from significance test of effect estimate from original study,
% % pr: Two-sided p-value from significance test of effect estimate from replication study,
% % po1: One-sided p-value from significance test of effect estimate from original study (in the direction of the original effect estimate),  
% % pr1: One-sided p-value from significance test of effect estimate from replication study (in the direction of the original effect estimate),
% % pm_belief: Peer belief about whether replication effect estimate will achieve statistical significance elicited through prediction market (only available for EERP and SSRP),
% % no: Sample size in original study,
% % nr: Sample size in replication study,

% % Goal: In which domains, a significant proportion of replication studies were conducted in a different country or language setting compared to the original study?

% % 
% % """
% %         \end{minted}

% % Example 3: meta\_regression\_raw||2||0

% % \begin{minted}[frame=lines,
% % baselinestretch=1,
% % bgcolor=Box3Color,
% % fontsize=\footnotesize, breaklines, breakanywhere, breaksymbolleft={}, breaksymbolright={}]{python}
% % task =
% % f"""\

% % Load all datasets using python using provided paths.
% % Paths: ['../DiscoveryBench/meta_regression_raw/meta-regression_replication_success_data_heterogeneity_in_replication_projects.csv', '../DiscoveryBench/meta_regression_raw/meta-regression_study_data_heterogeneity_in_replication_projects.csv']

% % Dataset name: meta-regression_study_data_heterogeneity_in_replication_projects.csv
% % Dataset description: Dataset contains information about original & replication studies. Original & replication specific columns may be appended by o & r
% % Brief description of columns:
% % id: Unique id for each O/R pair,
% % title: Title of the research study,
% % authors.o: Names of Original paper's authors,
% % pub_year: Year of Publication of the study,
% % journal: Journal in which the study was published,
% % volume: Volume Number of the journal,
% % issue: Issue Number of the journal,
% % discipline: Discipline of original paper. One of the following: Social, Cognitive or Economics,
% % length: Number of pages of original paper,
% % citations: Number of citations of original paper,
% % effect_size.o: Standardized effect size of original paper,
% % p_value.o: P-value of original paper,
% % n.o: Sample size of original paper,
% % effect_type: Type of effect tested. One of the following: main effect, correlation, interaction,
% % effect_size.r: Standardized effect size of replication,
% % p_value.r: P-value of replication,
% % n_planned.r: Planned sample size of replication,
% % n.r: Sample size of replication,
% % power.o: Post hoc power based on original effect size,
% % power.r: Post hoc power based on replication effect size,
% % power_planned.r: Planned power of the replication based on planned N and original ES,
% % experiment_country.o: Country where original experiment was conducted,
% % experiment_country.r: Country where replication is to be conducted,
% % experiment_language.o: Language used with subjects in original experiment (English, German, Dutch, Polish, Hebrew, French, Italian, Arabic, Spanish, Korean),
% % experiment_language.r: Language to be used with subjects in replication (English, Polish, German, Dutch, Italian, Portuguese, Malay, Turkish, Czech, Arabic, Spanish),
% % online.o: If the original experiment was conducted online (1: yes, 0: no),
% % online.r: If the replication was conducted online (1: yes, 0: no),
% % compensation.o: Compensation in original experiment (credit, cash, nothing, mixed),
% % compensation.r: Compensation in replication (credit, cash, nothing, mixed),
% % subjects.o: Type of subjects used in original experiment (students, online, anyone, community),
% % subjects.r: Type of subjects used in replication (students, online, anyone, community),
% % endprice: Final market price in prediction market,
% % transactions: Number of transactions in prediction market,
% % trading_volume: Total volume of traded stocks in prediction market,
% % replicated: Binary outcome variable; study is replicated if p <= 0.05 and effect goes in the same direction as the original,
% % project: The replication project that the study was in (ml1: Many Labs 1, ml3: Many Labs 3, rpp: Psychology, ee: Experimental Economics),
% % relative_es: The continuous outcome variable; the standardized replication effect size to the original effect | relative effect size = (replication effect size / original effect size),
% % n_authors.o: Number of authors in original study,
% % n_authors.r: Number of authors in replication,
% % author_citations_avg.o: Average number of citations of authors in original study,
% % author_citations_max.o: Number of citations of the author in original study with the highest citation count,
% % authors_male.o: Ratio of male authors in original study,
% % seniority.o: Most senior author in the original paper (Professor, Associate Professor, Assistant, Researcher, Assistant Professor),
% % author_citations_avg.r: Average number of citations of authors in replication study,
% % author_citations_max.r: Number of citations of the author in original study with the highest citation count,
% % authors_male.r: Ratio of male authors in replication,
% % seniority.r: Most senior author in the original paper (Professor, Associate Professor, Assistant, Researcher, Assistant Professor),
% % aggregated: Aggregated column,
% % lab_id: Unique id for each replication lab,
% % es_80power: Standardized effect size required in replication to achieve 80% power,
% % same_country: Original study and replication are in the same country,
% % same_language: Original study and replication are in the same language,
% % same_online: Original study and replication are both conducted online,
% % same_subjects: Original study and replication use same type of subjects,
% % us_lab.o: Original experiment lab in the US,
% % us_lab.r: Replication experiment lab in the US,
% % drop: Drop column,
% % Dataset name: meta-regression_replication_success_data_heterogeneity_in_replication_projects.csv
% % Dataset description: Data from four large-scale replication projects
% % Brief description of columns:
% % study: Study identifier, usually names of authors from original study,
% % project: Name of replication project,
% % ro: Effect estimate of original study on correlation scale,
% % rr: Effect estimate of replication study on correlation scale,
% % fiso: Effect estimate of original study transformed to Fisher-z scale,
% % fisr: Effect estimate of replication study transformed to Fisher-z scale,
% % se_fiso: Standard error of Fisher-z transformed effect estimate of original study,
% % se_fisr: Standard error of Fisher-z transformed effect estimate of replication study,
% % po: Two-sided p-value from significance test of effect estimate from original study,
% % pr: Two-sided p-value from significance test of effect estimate from replication study,
% % po1: One-sided p-value from significance test of effect estimate from original study (in the direction of the original effect estimate),
% % pr1: One-sided p-value from significance test of effect estimate from replication study (in the direction of the original effect estimate),
% % pm_belief: Peer belief about whether replication effect estimate will achieve statistical significance elicited through prediction market (only available for EERP and SSRP),
% % no: Sample size in original study,
% % nr: Sample size in replication study,

% % Goal: Which domain tend to have longer original papers?

% % 

% % """
% %         \end{minted}

% \subsection{Humanities}

% archaeology $\rightarrow$ 34 $\rightarrow$ 0 with \emph{domain knowledge}

% \begin{minted}[frame=lines,
% baselinestretch=1,
% bgcolor=Box3Color,
% fontsize=\footnotesize, breaklines, breakanywhere, breaksymbolleft={}, breaksymbolright={}]{python}
% task =
% f"""\

% Load all datasets using python using provided paths. Paths: ['../DiscoveryBench/archaeology/time_series_data.csv', '../DiscoveryBench/archaeology/capital.csv', '../DiscoveryBench/archaeology/pollen_openness_score_Belau_Woserin_Feeser_et_al_2019.csv']. 

% Dataset name: time_series_data.csv
% Dataset description: This dataset provides a detailed quantification of archaeological findings over various time periods, measured in Z values for different cultural and economic indicators such as tools, house sizes, materials, and monument data. 
% Brief description of columns: 
% CE: Common Era (BCE x (-1))
% calBP: Calibrated years before the present
% kde_all_mean: Mean of kernel density estimation of all data points
% kde_all_std: Standard deviation of kernel density estimation of all data points
% kde_all_detrend: KDE of data points after detrending
% g_all_mean: Mean of KDE growth rates
% g_all_std: Standard deviation of KDE growth rates
% pollen: Pollen data of Belau Lake
% pollen_inter: Interpolated and forward filled missing pollen values
% pollen_detrend: Detrended pollen values from interpolated pollen values
% pollen_inter_100: Rolling mean of the interpolated pollen data with a window size of 100
% pollen_grate_100: Percentage change of interpolated pollen data
% HatchetSword: Z values for Hatchets and Swords
% HatchetSword_inter: Interpolated z values for Hatchets and Swords
% Dagger: Z values for Daggers
% Dagger_inter: Interpolated z value for Daggers
% HouseSize: Z values for House Size in meter squared
% HouseSize_inter: Interpolated z values for House Sizes in meter squared
% CopperGold: Z values for Copper and Gold
% CopperGold_inter: Interpolated z values for Copper and Gold
% Amber: Z values for Amber
% Amber_inter: Interpolated z values for Amber
% MonumentCount: Z values for Monument Count
% MonumentCount_inter: Interpolated z values for Monument Count
% Depot: Z values for Depot
% Depot_inter: Interpolated z values for Depot
% Sickle: Z values for Sickle
% Sickle_inter: Interpolated z values for Sickle
% AxesCelts: Z values for Axes and Celts
% AxesCelts_inter: Interpolated z values for Axes and Celts
% MonumentSize: Z values for Monument Size
% MonumentSize_inter: Interpolated z values for Monument Size
% PotteryForm: Z values for Pottery Form
% PotteryForm_inter: Interpolated z values for Pottery Form
% PotteryDecoration: Z values for Pottery Decoration
% PotteryDecoration_inter: Interpolated z values for Pottery Decoration

% Dataset name: capital.csv
% Dataset description: This dataset contains archaeological data of various forms of capital across different prehistoric periods.
% Brief description of columns: 
% BCE: Before Common Era
% ZAxtSchwert: Z values for Hatchets and Swords
% ZDolch: Z values for Daggers
% Zhausgr: Z values for House Size
% ZCU_AU: Z values for Copper and Gold
% Zamber: Z values for Amber
% ZMonument: Z values for Monument Count
% ZHort: Z values for Depot
% ZSichel: Z values for Sickle
% ZBeil: Z values for Axes and Celts
% ZMW: Z values for Monument Size
% ZKeform: Z values for Pottery Form
% Zkeverz: Z values for Pottery Decoration

% Dataset name: pollen_openness_score_Belau_Woserin_Feeser_et_al_2019.csv
% Dataset description: Records of pollen data's PCA & interpolations acrosss sites.
% Brief description of columns: 
% calBP: Calibrated years Before Present (1950 AD)
% CE: Common Era
% Belau_PC1: PC1 of principal components for pollen in Belau
% Woserin_PC1: PC1 of principal components for pollen in Woserin
% Belau_PC1_inter: Interpolated PC1 for the Belau site
% Woserin_PC1_inter: Interpolated PC1 for the Woserin site
% MEAN: The average of the interpolated PC1 for the Belau and Woserin sites
% SMOOTH_MEAN_50y: Smoothed averages of the PC1 over 50 years
% SMOOTH_MEAN_100y: Smoothed averages of the PC1 over 100 years
% SMOOTH_MEAN_150y: Smoothed averages of the PC1 over 150 years
% SMOOTH_MEAN_200y: Smoothed averages of the PC1 over 200 years
% SMOOTH_MEAN_250y: Smoothed averages of the PC1 over 250 years 

% Goal: In what centuries did we see a steep dip in growth which rises to attain the highest peak of the past 500 years around 1400 BCE? Additionally, we provide some hints that might be useful to solve the task. 

% Domain Knowledge: 
% 1. Symbolic capital consists of Hatchet & Swords, Daggers, House Size. 
% 2. Social Capital consists of Copper and Gold, Amber, Monument Count 
% 3. Cultural Capital consists of Diversity of Pottery form, Diversity of Pottery Decoration. 
% 4. Economic Capital consists of Depot, Sickle, Axes & Celts, Monument Size. 
% 5. Human impact or landscape openness, respectively, as reflected in the pollen data, can be used as a demographic indicator based on the assumption that an increasing population density leads to increasing woodland clearance due to an increasing demand for resources including wood, agricultural land and settlement areas. Each sample from the pollen record used in the principal component analysis is absolutely dated and therefore the openness score (PC 1.) can be plotted as a time series, expressing human induced land clearance. The 'Belau_PC1' of pollen data (pollen_openness_score_Belau_Woserin_Feeser_et_al_2019.csv) has been assumed to reflect openness. Original openness score through 'Belau_PC1', 100-year smoothed openness score, and linear interpolation of openness score have been used to signify growth. The original openness score (Belau_PC1), the 100-year smoothed openness score (pollen_inter_100), and the linear interpolation of the openness score (pollen_inter) have been used to signify openness.
% 6. Demographic growth manifests itself in growth set bringing with it a further opening of the landscape. Opening of landscape corresponds to higher growth rates. During the Early Neolithic, we are dealing with a population growth that goes hand in hand with the opening up of vegetation and the cultivation of the landscape. The growth rate is defined as the percentage change of the 100-year smoothed openness score (pollen_grate_100).
% 7. Time series analysis and PCA are done in 100 year bins.

% """
%         \end{minted}
